\section[Riepilogo del documento]{Riepilogo del documento}

{\small\arrayrulecolor{rulec}\begin{xltabular}{\linewidth}{>{\hsize=0.4348\hsize\raggedright\arraybackslash}X>{\hsize=1.5652\hsize\raggedright\arraybackslash}X}
\toprule

\textbf{Documento} & \textbf{symphony\_\allowbreak{}version\_\allowbreak{}log} (unico file, aggiornato a ogni revisione). Fino alla rev. 1.5 si chiamava \code{modify\_\allowbreak{}request1910}: tutto il suo contenuto è riportato qui, invariato \\
\midrule[0.15pt]

\textbf{Revisione corrente} & \textbf{1.40} del 03/11/2026 \textendash{} R-21: stock dei componenti e sostituzione di C217/C218\newline{}1.39 del 03/11/2026 \textendash{} R-20: margine di oscillazione del quarzo Y1, chiuso senza modifiche\newline{}1.38 del 03/11/2026 \textendash{} R-17: pin SHDN del TCAN3414 collegato a GND\newline{}1.37 del 03/11/2026 \textendash{} R-18: precisione della catena PT1000, chiuso senza modifiche\newline{}(le revisioni 1.37--1.40 sono descritte nelle ultime quattro entry del log; da qui in poi elenco della versione HTML:)\newline{}\textbf{1.36} del 02/\allowbreak{}11/\allowbreak{}2026 – R-14 chiuso senza modifiche: C39 resta 10 µF 0402 10 V \\
\midrule[0.15pt]

\textbf{1.35 del 02/\allowbreak{}11/\allowbreak{}2026 – R-13 chiuso senza modifiche: pull-up I2C1 tramite JP201 aperto, standard di progetto} & 1.34 del 02/\allowbreak{}11/\allowbreak{}2026 – R-19: supervisore IC205 escludibile (R204 0R montata, R205 0R DNP verso PC4); misura M-03 e note firmware F-05/\allowbreak{}F-10 \\
\midrule[0.15pt]

\textbf{1.33 del 02/\allowbreak{}11/\allowbreak{}2026 – R-16: ST del TPS2121 collegato a GND (Brik81); OV2 lasciato a GND per decisione} & 1.32 del 02/\allowbreak{}11/\allowbreak{}2026 – R-07 chiuso senza modifiche: il TCXO Y201 resta ECS-TXO-20CSMV4-320-AY-TR (−30…+85 °C) \\
\midrule[0.15pt]

\textbf{1.31 del 02/\allowbreak{}11/\allowbreak{}2026 – R-12: +12V, +3V3\_\allowbreak{}CONN e RS485 di J708 non usati per scelta, nota sullo schema IF} & 1.30 del 02/\allowbreak{}11/\allowbreak{}2026 – R-06/\allowbreak{}R-24 chiusi senza modifiche: il backup dell'RTC serve solo per pochi secondi \\
\midrule[0.15pt]

\textbf{1.29 del 02/\allowbreak{}11/\allowbreak{}2026 – R-04 chiuso senza modifiche allo schema: limite di potenza e sequenza TX affidati al firmware; raccolta delle note per il firmware (§41)} & 1.28 del 30/\allowbreak{}10/\allowbreak{}2026 – R-10 chiuso: 8 condensatori RF da 100/\allowbreak{}56 pF passano a Murata GQM1875 0603; R-09 chiuso con J15 lasciato com'è \\
\midrule[0.15pt]

\textbf{1.27 del 30/\allowbreak{}10/\allowbreak{}2026 – Gruppo B: test point SMD (R-08), C261/\allowbreak{}C702/\allowbreak{}C704 Murata GJM (R-10, parte), pin schermo USB-C S1→SH (R-09, parte)} & 1.26 del 30/\allowbreak{}10/\allowbreak{}2026 – R-05: codici JLC di R206/\allowbreak{}R207, R202 e C35 corretti; R-22: datasheet di U701 (74AUP2G14 Diodes) \\
\midrule[0.15pt]

\textbf{1.25 del 30/\allowbreak{}10/\allowbreak{}2026 – R-03/\allowbreak{}R-11: L242 torna a 1 µH (Taiyo Yuden NRS2012T1R0NGJV) e C260 a 27 pF, come la EVB-02 Qorvo} & 1.24 del 27/\allowbreak{}10/\allowbreak{}2026 – Revisione indipendente R-01/\allowbreak{}R-02: codici JLC di U301 (STM32L562RET6P) e IC1/\allowbreak{}IC5 (TLV9051IDBVR) corretti \\
\midrule[0.15pt]

\textbf{1.23 del 27/\allowbreak{}10/\allowbreak{}2026 – S-11: D202 sull'I²C1 sostituito con un TVS a 5 V (PESD5V0L2BT)} & 1.22 del 27/\allowbreak{}10/\allowbreak{}2026 – tolta l'etichetta UHF.EN\_\allowbreak{}RX\_\allowbreak{}LNA senza collegamenti nel foglio MCU \\
\midrule[0.15pt]

\textbf{1.21 del 24/\allowbreak{}10/\allowbreak{}2026 – origine del footprint CPH3225A al centro; testo spurio tolto dal foglio principale; U701 verificato} & 1.20 del 24/\allowbreak{}10/\allowbreak{}2026 – S-09 R410 DNP; simbolo di non collegato su tutti i GPIO liberi dell'MCU; VBAT del DS3231M ricollegato da Brik81 \\
\midrule[0.15pt]

\textbf{1.19 del 24/\allowbreak{}10/\allowbreak{}2026 – S-07 R245 in 0603; S-14 C307 in 0603 e L241 in 0402; esito di S-01, S-08, S-10…S-13} & 1.18 del 24/\allowbreak{}10/\allowbreak{}2026 – S-06: codici delle resistenze di precisione e PT1000 0603 da acquistare a parte; S-06 chiuso \\
\midrule[0.15pt]

\textbf{1.17 del 24/\allowbreak{}10/\allowbreak{}2026 – S-06: uscite dei sensori PT1000 collegate all'ADC (V\_\allowbreak{}TEMP\_\allowbreak{}PA su PB1, V\_\allowbreak{}TEMP\_\allowbreak{}BUCK su PC1)} & 1.16 del 24/\allowbreak{}10/\allowbreak{}2026 – S-05: condensatori 10 µF + 100 nF sull'ingresso IN1 (5 V dock) del TPS2121; S-05 chiuso \\
\midrule[0.15pt]

\textbf{1.15 del 24/\allowbreak{}10/\allowbreak{}2026 – S-04: valore del quarzo Y1 (24 MHz, CL 10 pF) e codice JLCPCB di C1/\allowbreak{}C2 corretti; S-04 chiuso} & 1.14 del 24/\allowbreak{}10/\allowbreak{}2026 – S-03 verificato e chiuso \\
\midrule[0.15pt]

\textbf{1.13 del 24/\allowbreak{}10/\allowbreak{}2026 – nomi alle reti RF e footprint standard sui componenti del percorso RF} & 1.12 del 24/\allowbreak{}10/\allowbreak{}2026 – netclass RF, coppia differenziale RX del Si4464 e guida al layout della parte RF \\
\midrule[0.15pt]

\textbf{1.11 del 24/\allowbreak{}10/\allowbreak{}2026 – induttore di alimentazione dello stadio finale del PA (L242) assegnato: Murata RF 22 nH 1.9 A, con C260 corretto a 33 pF} & 1.10 del 24/\allowbreak{}10/\allowbreak{}2026 – codice JLCPCB unificato nel campo \textquotedblleft{}JLCPCB Parts\textquotedblright{} su tutto il progetto \\
\midrule[0.15pt]

\textbf{1.9 del 23/\allowbreak{}10/\allowbreak{}2026 – rete RX del Si4464: balun a traliccio sostituito dalla rete AN643 a 4 elementi (S-03)} & 1.8 del 23/\allowbreak{}10/\allowbreak{}2026 – U402 (USB-UART) spostato da UART4 (PA0/\allowbreak{}PA1) a USART1 (PA9/\allowbreak{}PA10) per il bootloader seriale \\
\midrule[0.15pt]

\textbf{1.7 del 23/\allowbreak{}10/\allowbreak{}2026 – S-02 rifatto in SMD: 10 × Seiko CPH3225A in parallelo al posto del supercap passante della 1.6} & 1.6 del 21/\allowbreak{}10/\allowbreak{}2026 – supercap passante (THT) sul VBAT dell'RTC, \textbf{respinta} da Brik81 il 21/\allowbreak{}10 perché non SMD \\
\midrule[0.15pt]

\textbf{Revisioni precedenti} & 0.1 · 1.0 · 1.1 · 1.2 · 1.3 del 19/\allowbreak{}10/\allowbreak{}2026 (Parte B) · 1.4 · 1.5 del 20/\allowbreak{}10/\allowbreak{}2026 (Parti C e D), contenuto invariato \\
\midrule[0.15pt]

\textbf{Progetto} & SYMPHONY\_\allowbreak{}V2, branch SYM\_\allowbreak{}V2, KiCad 10.0 \\
\midrule[0.15pt]

\textbf{Fogli} & SYMPHONY (radice), POWER, IF, MCU, UHF, SENSORS, MOUNTING\_\allowbreak{}HOLES \\
\midrule[0.15pt]

\textbf{Redatto da} & Claude Code (Claude Opus 5.5) su richiesta di Brik81 \\
\midrule[0.15pt]

\bottomrule
\end{xltabular}}

\subsection[File di progetto modificati per revisione]{File di progetto modificati per revisione}

{\small Testo ripreso dalla voce dell'indice HTML relativa alla Parte M (nel sorgente HTML finito in quella voce per un errore di sostituzione; riportato per intero):}

Rev. 1.36: nessun file di progetto modificato. Rev. 1.35: nessun file di progetto modificato. Rev. 1.34: nessun file di progetto modificato da Claude Code (R204/\allowbreak{}R205 aggiunte da Brik81 in UHF). Rev. 1.33: nessun file di progetto modificato da Claude Code (ST collegato da Brik81 in POWER). Rev. 1.32: nessun file di progetto modificato. Rev. 1.31: \code{IF.\allowbreak{}kicad\_\allowbreak{}sch}, solo una nota di testo (copia precedente in \code{report\_\allowbreak{}src/\allowbreak{}bak\_\allowbreak{}pre131/\allowbreak{}}). Rev. 1.30: nessun file di progetto modificato. Rev. 1.29: nessun file di progetto modificato. Rev. 1.28: \code{UHF.\allowbreak{}kicad\_\allowbreak{}sch} (copia precedente in \code{report\_\allowbreak{}src/\allowbreak{}bak\_\allowbreak{}pre128/\allowbreak{}}). Rev. 1.27: \code{UHF.\allowbreak{}kicad\_\allowbreak{}sch} e \code{POWER.\allowbreak{}kicad\_\allowbreak{}sch} (copie precedenti in \code{report\_\allowbreak{}src/\allowbreak{}bak\_\allowbreak{}pre127/\allowbreak{}}). Rev. 1.26: \code{UHF.\allowbreak{}kicad\_\allowbreak{}sch} e \code{UHF\_\allowbreak{}LIB/\allowbreak{}SYMBOL/\allowbreak{}IC.\allowbreak{}kicad\_\allowbreak{}sym}, solo il campo Datasheet di 74AUP2G14 (copie precedenti in \code{report\_\allowbreak{}src/\allowbreak{}bak\_\allowbreak{}pre126/\allowbreak{}}). Rev. 1.25: \code{UHF.\allowbreak{}kicad\_\allowbreak{}sch} (copia precedente in \code{report\_\allowbreak{}src/\allowbreak{}bak\_\allowbreak{}pre125/\allowbreak{}}). Rev. 1.24: nessun file di progetto modificato da Claude Code (le modifiche le ha fatte Brik81 in KiCad); netlist prima e dopo nella cartella temporanea della sessione. Rev. 1.23: \code{IF.\allowbreak{}kicad\_\allowbreak{}sch} prima della modifica salvato in \code{report\_\allowbreak{}src/\allowbreak{}bak\_\allowbreak{}pre123/\allowbreak{}}. Rev. 1.21: \code{CPH3225A.\allowbreak{}kicad\_\allowbreak{}mod} prima della modifica salvato in \code{bak\_\allowbreak{}pre121/\allowbreak{}}. Rev. 1.20: \code{MCU.\allowbreak{}kicad\_\allowbreak{}sch} prima della modifica salvato in \code{bak\_\allowbreak{}pre120/\allowbreak{}}. Rev. 1.19: \code{UHF.\allowbreak{}kicad\_\allowbreak{}sch}, \code{MCU.\allowbreak{}kicad\_\allowbreak{}sch} prima della modifica salvati in \code{bak\_\allowbreak{}pre119/\allowbreak{}}. Rev. 1.18: \code{CONNECTORS.\allowbreak{}kicad\_\allowbreak{}sch} prima della modifica salvato in \code{bak\_\allowbreak{}pre118/\allowbreak{}}. Rev. 1.17: \code{CONNECTORS.\allowbreak{}kicad\_\allowbreak{}sch}, \code{MCU.\allowbreak{}kicad\_\allowbreak{}sch}, \code{SYMPHONY.\allowbreak{}kicad\_\allowbreak{}sch} prima della modifica salvati in \code{bak\_\allowbreak{}pre117/\allowbreak{}}. Rev. 1.16: \code{POWER.\allowbreak{}kicad\_\allowbreak{}sch} prima della modifica salvato in \code{bak\_\allowbreak{}pre116/\allowbreak{}}. Rev. 1.15: \code{MCU.\allowbreak{}kicad\_\allowbreak{}sch} prima della modifica salvato in \code{bak\_\allowbreak{}pre115/\allowbreak{}}. Rev. 1.14: chiusura di S-03

\clearpage
\section[Parte A – Registri]{Parte A – Registri}

\subsection[Registro revisioni]{Registro revisioni}

\begin{notebox}
\textbf{Regola di approvazione.} Nessuna modifica ai file di progetto viene eseguita senza il consenso esplicito di una persona. Ogni revisione e ogni singola modifica riporta chi l'ha approvata, quando, e chi l'ha eseguita. Le analisi e le proposte non approvate restano nella sezione \textquotedblleft{}Proposte in attesa\textquotedblright{} e non toccano i file. Orari in ora locale (CEST, UTC+2), presi dal log della sessione.
\end{notebox}

{\scriptsize\arrayrulecolor{rulec}\setlength{\tabcolsep}{4pt}\renewcommand{\arraystretch}{1.1}
\begin{xltabular}{\linewidth}{>{\hsize=0.2575\hsize\raggedright\arraybackslash}X>{\hsize=1.0522\hsize\raggedright\arraybackslash}X>{\hsize=1.7116\hsize\raggedright\arraybackslash}X>{\hsize=0.9020\hsize\raggedright\arraybackslash}X>{\hsize=0.5995\hsize\raggedright\arraybackslash}X>{\hsize=1.4773\hsize\raggedright\arraybackslash}X}
\toprule
\textbf{Rev.} & \textbf{Data e ora} & \textbf{Descrizione} & \textbf{Eseguito da} & \textbf{Approvato da (persona)} & \textbf{Testo dell'approvazione} \\
\midrule
\endhead
\midrule[0.15pt]\multicolumn{6}{r}{\footnotesize\itshape continua sulla pagina seguente\ldots}\\\bottomrule\endfoot
\bottomrule\endlastfoot
0.1 & 19/\allowbreak{}10/\allowbreak{}2026 12:01 richiesta\newline{}12:01–12:30 analisi & Revisione del foglio UHF: solo analisi, nessuna modifica ai file & Claude Code & – (nessuna modifica) & Brik81, 12:01: \textquotedblleft{}…prima di modificare nel caso avvertimi\textquotedblright{} \\\midrule[0.15pt]
1.0 & approvata 12:34\newline{}eseguita 12:35–12:46\newline{}report 12:48 & TCXO su XIN, match TX del Si4464, APC, modifiche minori (esclusa SDN), via termiche IC2, creazione di questo report & Claude Code & \textbf{Brik81} & 12:34: \textquotedblleft{}2) … provvedi, abilitato a modificare · 3) se sei sicuro al 100\% provvedi, abilitato a modificare · 5) ok mi fido, provvedi · modifiche minori: falle tutte tranne l'ultima SDN del Si4464\textquotedblright{} \\\midrule[0.15pt]
1.1 & approvata 15:49\newline{}eseguita 15:51–15:53\newline{}report 15:55 & R29/\allowbreak{}R31 scambiati (POWER), simbolo P-MOS per Q701, VDD dell'HMC349 su 5V0. Aggiunti registro revisioni, verifica passivi di IC2 (§9) e verifica codici JLCPCB (§10) & Claude Code & \textbf{Brik81} & 15:49: \textquotedblleft{}Ok procedi con le modifiche rimaste da fare verificate, le approvo\textquotedblright{} \\\midrule[0.15pt]
1.2 & approvata 17:04\newline{}eseguita 17:05\newline{}report 17:06 & Proposta P3 (§9): rimossi R244 e C253. Registrata la decisione D1 sulla priorità del power mux & Claude Code & \textbf{Brik81} & 17:04: \textquotedblleft{}per quanto riguarda le modifiche da fare dopo la rev v1.1 il punto 9, scegli tu chi rimuovere\textquotedblright{} \\\midrule[0.15pt]
1.3 & approvata 17:11\newline{}eseguita 17:13\newline{}report 17:15 & Proposta P4 (§10): codici JLCPCB corretti e assegnati (condensatori e induttori Murata verificati sul catalogo JLCPCB), cambio footprint di L243 e L703. Regola tensione aggiornata a ≥ 3× & Claude Code & \textbf{Brik81} & 17:11: \textquotedblleft{}ok parliamo del punto 10 … cmq modifichiamo, controllare che i componenti siano murata da jlcpcb … verificando il codice JLCPCB, come power rating direi che deve stare tra 3 e 5 volte l'alimentazione\textquotedblright{} \\\midrule[0.15pt]
\rowcolor{extrow}1.3+ & 19/\allowbreak{}10/\allowbreak{}2026 tra 17:13 e 17:26 (rilevata il 20/\allowbreak{}10 09:22) & Rimossi C226, C227, C228 (UHF) e J3, R10, R14, R15 (IF, connettore di debug vecchio, già DNP). Rilevata confrontando la netlist con quella di fine sessione & In KiCad, fuori dalla sessione Claude Code (autore non registrato, presumibilmente Brik81: \textbf{da confermare}) & – (non registrato) & – \\\midrule[0.15pt]
1.4 & richiesta 20/\allowbreak{}10/\allowbreak{}2026 09:19\newline{}analisi 09:19–09:31\newline{}report 09:31 & Verifica completa di tutti i fogli (Parte C, §13–§15) e nuova impaginazione del documento. Storia delle rev. 0.1–1.3 invariata. \textbf{Nessuna modifica ai file di progetto}; nuove proposte P8–P17 & Claude Code & – (nessuna modifica) & Brik81, 20/\allowbreak{}10: \textquotedblleft{}mi fai un controllo di tutto il progetto? foglio per foglio … una parte finale in cui confermi che ogni foglio abbia i vari componenti dimensionati e connessi bene, non ti preoccupare del codice jlcpcb\textquotedblright{} \\\midrule[0.15pt]
\rowcolor{extrow}1.4+ & 20/\allowbreak{}10/\allowbreak{}2026 tra 09:31 e 14:38 (salvataggio KiCad) & Correzioni P8–P12 e altre (Parte D, §16): IC501 non più a specchio, VDD12 NC, C1/\allowbreak{}C2 10 pF, R202 0 Ω /\allowbreak{} R203 DNP, U701 pin 3 a GND, switch RF1↔RF2, R206/\allowbreak{}R207 4k7, SENSORS completato con il secondo canale & Brik81, in KiCad (fuori sessione Claude Code) & \textbf{Brik81} (autore) & Brik81, chat 20/\allowbreak{}10: \textquotedblleft{}IF: … è stato specchiato per errore, risolto · MCU: … cavando la connessione vdd12 e lasciandola NC e cambiato i valori del C del clock\textquotedblright{}; RF1/\allowbreak{}RF2 scambiati dopo la discussione sullo switch \\\midrule[0.15pt]
1.5 & approvata 20/\allowbreak{}10/\allowbreak{}2026 17:02\newline{}eseguita 17:07–17:12\newline{}report 17:15 & S-01 (sch\_\allowbreak{}check\_\allowbreak{}2010): supervisore TPS3825-33 IC205 sulla SDN del Si4464, \textasciitilde{}MR su PC4 con pull-up, LED rosso di stato. Registrate le modifiche 1.4+ & Claude Code & \textbf{Brik81} & 17:02: \textquotedblleft{}ok allora fai te per me, collega mr con un pullup fatto bene al gpio che ritieni opportuno … cercami sia la resistenza che led su jlcpcb, piazzameli con il footprint corretto e con il jlcpcb part … per il cap … usa quello [100n 0402] … e il resto dei collegamenti falli tu ok? … vai parti\textquotedblright{} \\\midrule[0.15pt]
1.6 & approvata 21/\allowbreak{}10/\allowbreak{}2026 tra 17:44 e 17:55\newline{}eseguita 17:55–17:58\newline{}report 18:00 & S-02 (sch\_\allowbreak{}check\_\allowbreak{}2010): BT401 (MS621FE, non acquistabile dall'Italia e senza ricarica) sostituita da C415 supercondensatore 0.33 F 5.5 V, caricato dal 3V3 tramite D401 BAT54WS e R412 1 kΩ. Documento rinominato in symphony\_\allowbreak{}version\_\allowbreak{}log & Claude Code & \textbf{Brik81} & \textquotedblleft{}si procedi con la modifica di S-02 RTC, pero nota che le alimentazioni sono tutte a 3v3 … volgio la modifica del report assolutamente, cambia il nome del documento in \textquotedblright{}symphony\_\allowbreak{}version\_\allowbreak{}log.pdf\textquotedblleft{} … rpocedi pure\textquotedblright{}. Prima: \textquotedblleft{}dai fai cosi, allora modifica eliminando la battteria e mettendo il supercap … che duri almeno 15 ore\textquotedblright{} \\\midrule[0.15pt]
1.7 & richiesta 23/\allowbreak{}10/\allowbreak{}2026\newline{}eseguita 17:18 & S-02 rifatto in SMD. La rev. 1.6 (C415 0.33 F passante) era stata respinta il 21/\allowbreak{}10: \textquotedblleft{}deve essere SMD!!!! OBBLIFATORIO NIENTE FORO PASSANTE\textquotedblright{}. C415…C424 = 10 × Seiko CPH3225A 11 mF 3.3 V in parallelo su RTC\_\allowbreak{}VBAT; D401 e R412 invariati & Claude Code & \textbf{Brik81} & \textquotedblleft{}nel foglio MCU ho incollato il condenstore da usare e il diodo D401 e R412 come avevi fatto tu, è da finire il collegamento, metti in parallelo N condensatori … direi una aujtonomia di 5H è sufficente … finsici tu la connessione verso Vbat\textquotedblright{} \\\midrule[0.15pt]
1.8 & proposta 21/\allowbreak{}10/\allowbreak{}2026, approvata 23/\allowbreak{}10/\allowbreak{}2026\newline{}MCU eseguita da Brik81; collegamento tra fogli eseguito da Claude Code il 23/\allowbreak{}10 & U402 CH340 spostato da UART4 (PA0/\allowbreak{}PA1), che il bootloader di sistema STM32L562 non supporta (AN2606), a USART1 (PA9/\allowbreak{}PA10), che lo supporta. Tutti i nomi passano a UART1 & Brik81 (MCU), Claude Code (fogli principale e IF) & \textbf{Brik81} & \textquotedblleft{}la uart dovrei aver gia fatto io verifica\textquotedblright{} → \textquotedblleft{}si vai, proceid con la modifica del nome con UART1\textquotedblright{} \\\midrule[0.15pt]
1.9 & proposta e approvata 23/\allowbreak{}10/\allowbreak{}2026\newline{}eseguita da Claude Code il 23/\allowbreak{}10 & S-03 (sch\_\allowbreak{}check\_\allowbreak{}2010): il balun RX a traliccio (L205 18n, C219 7.5p, C220 7.5p, L206 22n, C222/\allowbreak{}C223 100p) era dimensionato per circa 50 Ω differenziali, mentre l'ingresso del Si4464 a 434 MHz è 480 Ω ∥ 0.97 pF. Sostituito con la rete AN643 a 4 elementi (valori ottimizzati di Tab. 5): 2.2p /\allowbreak{} 56n /\allowbreak{} 4.7p /\allowbreak{} 62n & Claude Code & \textbf{Brik81} & \textquotedblleft{}andiamo a s-03\textquotedblright{} → \textquotedblleft{}ok procedi e speriamo bene\textquotedblright{}. 24/\allowbreak{}10: \textquotedblleft{}procedi con la correzione MA RICORDA, AGGIUNGI (O MODIFICA) NELLE PROPRIETA IL PARAMETRO JLCPCB PART DI OGNI COMPONENTE\textquotedblright{} \\\midrule[0.15pt]
1.10 & richiesta e approvata 24/\allowbreak{}10/\allowbreak{}2026\newline{}eseguita da Claude Code il 24/\allowbreak{}10 & Il codice di parte JLCPCB era scritto in 5 campi con nomi diversi (jlc, jlcpcb, JLCPCB Part, JLCPCB PART, JLCPCB), in alcuni componenti anche due volte. Unificato in un solo campo \textquotedblleft{}JLCPCB Parts\textquotedblright{} in tutti i fogli, nelle librerie simboli di progetto e nella lista campi BOM del progetto & Claude Code & \textbf{Brik81} & \textquotedblleft{}si rinomina tutti i campi li voglio unificati in \textquotedblright{}JLCPCB Parts\textquotedblleft{}\textquotedblright{} \\\midrule[0.15pt]
1.11 & proposta e approvata 24/\allowbreak{}10/\allowbreak{}2026\newline{}eseguita da Claude Code il 24/\allowbreak{}10 & L242, l'induttore di alimentazione dello stadio finale del QPA9510, non aveva codice JLCPCB (0.9 µH 1008). Nessun induttore Murata RF da circa 1 µH regge 1.7 A. Scelto 22 nH Murata LQW2BAN22NJ00L (1.9 A) come nella scheda di valutazione Qorvo; C260 portato da 27 a 33 pF per mantenere l'adattamento d'uscita. Verificati anche tutti gli altri induttori & Claude Code & \textbf{Brik81} & \textquotedblleft{}L242 SE PUOI ASSEGNARLO TU? STAI ATTENTO ALLE CORRENTI DI SATURAZIONE ECC ECC, E CHE SIA RF MURATA (MA NON SOLO PER LUI, PER TUTTI GLI INDUTTORI)\textquotedblright{} → \textquotedblleft{}OK PROCEDI\textquotedblright{} \\\midrule[0.15pt]
1.12 & richiesta e approvata 24/\allowbreak{}10/\allowbreak{}2026\newline{}eseguita da Claude Code il 24/\allowbreak{}10 & Aggiunte 6 netclass (RF\_\allowbreak{}DIFF, RF\_\allowbreak{}PA\_\allowbreak{}OUT, RF\_\allowbreak{}50R, RF\_\allowbreak{}CLK, PA\_\allowbreak{}BIAS, PWR) assegnate alle reti del foglio UHF. Le linee RXP/\allowbreak{}RXN del Si4464 diventano la coppia differenziale RF\_\allowbreak{}RX\_\allowbreak{}P /\allowbreak{} RF\_\allowbreak{}RX\_\allowbreak{}N. Guida al layout della parte RF (§24) & Claude Code & \textbf{Brik81} & \textquotedblleft{}sul foglio UHF modificami le netclass, rendi differenziali le linee, modifica le netclass e fammi un piccolo schema di dove e come andranno a livello di layout i vari componenti della parte rf, a cosa devo stare attento, alle dimensioni consigliate delle piste\textquotedblright{} \\\midrule[0.15pt]
1.13 & proposta 24/\allowbreak{}10 (§24.2, §24.5), approvata 24/\allowbreak{}10/\allowbreak{}2026\newline{}eseguita da Claude Code il 24/\allowbreak{}10 & Date le etichette a 28 reti RF che avevano nomi automatici, così le assegnazioni di netclass non si perdono se i componenti vengono rinumerati. 27 componenti del percorso RF passano da footprint HandSolder a footprint standard (pad più piccoli, meno capacità parassita) & Claude Code & \textbf{Brik81} & \textquotedblleft{}ok, modifica i pad con quelli handsolder dove serve, si dai un nome alle varie net\textquotedblright{} \\\midrule[0.15pt]
1.14 & verifica richiesta e approvata 24/\allowbreak{}10/\allowbreak{}2026\newline{}eseguita da Claude Code e Brik81 il 24/\allowbreak{}10 & Verifica finale di S-03 (rete RX AN643 del Si4464, rev. 1.9–1.13) e chiusura del punto. Nessuna modifica al progetto & Claude Code & \textbf{Brik81} & \textquotedblleft{}ALLORA controlliamo che sia tutto in regola con S-03 e chiudiamo s-03\textquotedblright{} \\\midrule[0.15pt]
1.15 & proposta e approvata 24/\allowbreak{}10/\allowbreak{}2026\newline{}eseguita da Claude Code il 24/\allowbreak{}10 & S-04 (sch\_\allowbreak{}check\_\allowbreak{}2010): il valore di Y1 era il nome del simbolo (\textquotedblleft{}Crystal\_\allowbreak{}GND24\textquotedblright{}) e non diceva frequenza e CL. Scritto \textquotedblleft{}24MHz 10pF\textquotedblright{}. Tolto uno spazio finale dal codice JLCPCB di C1 e C2. Circuito dell'oscillatore ricontrollato. S-04 chiuso & Claude Code & \textbf{Brik81} & \textquotedblleft{}sisi ok andiamo avanti\textquotedblright{} \\\midrule[0.15pt]
1.16 & proposta (S-05, P13 della rev. 1.4) e approvata 24/\allowbreak{}10/\allowbreak{}2026\newline{}eseguita da Claude Code il 24/\allowbreak{}10 & L'ingresso IN1 del TPS2121 (IC4), cioè il 5 V della dock, non aveva condensatori. Aggiunti C39 10 µF e C40 100 nF, entrambi 0402 per occupare il minimo spazio. S-05 chiuso & Claude Code & \textbf{Brik81} & \textquotedblleft{}si mettili, minimo spazio possibile e che siano su jlcpcb part\textquotedblright{} \\\midrule[0.15pt]
1.17 & proposta e approvata 24/\allowbreak{}10/\allowbreak{}2026\newline{}eseguita da Claude Code il 24/\allowbreak{}10 & S-06 (sch\_\allowbreak{}check\_\allowbreak{}2010): le uscite V\_\allowbreak{}TEMP\_\allowbreak{}PA e V\_\allowbreak{}TEMP\_\allowbreak{}BUCK del foglio SENSORS non arrivavano all'MCU. Collegate con etichette gerarchiche a due ingressi ADC liberi il più lontano possibile tra loro: V\_\allowbreak{}TEMP\_\allowbreak{}PA su PB1 (pin 26, ADC1\_\allowbreak{}IN16), V\_\allowbreak{}TEMP\_\allowbreak{}BUCK su PC1 (pin 9, ADC1\_\allowbreak{}IN2) & Claude Code & \textbf{Brik81} & \textquotedblleft{}falle opposte perche un sensore è nella parte del micro mentre l'altro sensore sarà nella parte del uhf\textquotedblright{} → \textquotedblleft{}si ok, si fai le etichette gerarchiche come le hai nominate tu che è come le voglio, come configurazione fai quella proposta da te\textquotedblright{} \\\midrule[0.15pt]
1.18 & proposta e approvata 24/\allowbreak{}10/\allowbreak{}2026\newline{}eseguita da Claude Code il 24/\allowbreak{}10 & Codici JLCPCB delle 8 resistenze di precisione del foglio SENSORS (0402, 0.1 \%, 25 ppm/\allowbreak{}°C, footprint invariato). PT1000 R\_\allowbreak{}RTD1/\allowbreak{}R\_\allowbreak{}RTD2: nessuna PT1000 0603 su JLCPCB; scelta Yageo Nexensos 32207638 (0603, classe B) da acquistare a parte e saldare a mano, footprint invariato. S-06 chiuso & Claude Code & \textbf{Brik81} & \textquotedblleft{}cercami i codici JLCPCB di R4–R10, R12 e delle PT1000. CHE RISPETTINO ASSOLUTAMENTE IL FOOTPRINT\textquotedblright{} → \textquotedblleft{}mi tocca fare A perché poi non ho mica spazio… indicami il codice mouser invece che jlcpcb e poi lo compro a parte e salderò a mano\textquotedblright{} \\\midrule[0.15pt]
1.19 & decisioni di Brik81 del 24/\allowbreak{}10/\allowbreak{}2026\newline{}eseguita da Claude Code il 24/\allowbreak{}10 & Rilievi minori di sch\_\allowbreak{}check\_\allowbreak{}2010: R245 passa a 0603 (S-07); C307 a 0603 e L241 a 0402 (S-14). Verificate le correzioni fatte da Brik81 (S-08 R203 DNP, S-10 etichetta 5v\_\allowbreak{}TX\_\allowbreak{}EN tolta). S-11, S-12 lasciati così per decisione, S-13 invariato (LED volutamente debole). S-01 risulta già risolto con IC205 & Claude Code & \textbf{Brik81} & \textquotedblleft{}s-07: ok sostituisci con 0603 … s-08: risolto … s-10: fixed … s-11: vanno bene così … s-12: lascia così … s-13: DEVE VEDERSI POCO … s-14: procedi come farebbe un senior hw designer\textquotedblright{} \\\midrule[0.15pt]
1.20 & proposta e approvata 24/\allowbreak{}10/\allowbreak{}2026\newline{}eseguita da Claude Code il 24/\allowbreak{}10 & S-09: R410 (pull-up dell'uscita 32KHZ del DS3231M) non montata. Simbolo di non collegato sui 22 GPIO liberi di U301 (PC3 compreso). Registrato anche il ricollegamento del pin VBAT del DS3231M fatto da Brik81 (salvataggio delle 17:50) & Claude Code & \textbf{Brik81} & \textquotedblleft{}ok si fai entrambe e metti NC in tutti gli altri pin non usati\textquotedblright{} \\\midrule[0.15pt]
1.21 & richiesta e approvata 24/\allowbreak{}10/\allowbreak{}2026\newline{}eseguita da Claude Code e Brik81 il 24/\allowbreak{}10 & Origine del footprint CPH3225A spostata al centro di pad e corpo (era fuori centro di 0.3 mm in X e Y). Il testo \textquotedblleft{}cd C:\textbackslash{}...\textquotedblright{} nel foglio principale è stato tolto da Brik81. U701 confermato come simbolo a più unità, non doppione & Claude Code & \textbf{Brik81} & \textquotedblleft{}u701 non è un doppione è un simbolo suddiviso in diversi simboli indipendenti, modifica l'origine, al 3D me la sbrigo io, cava la scritta cd ...\textquotedblright{} \\\midrule[0.15pt]
1.22 & proposta da Claude Code e approvata 27/\allowbreak{}10/\allowbreak{}2026\newline{}eseguita da Brik81 il 27/\allowbreak{}10 & Tolti nel foglio MCU l'ingresso dal bus UHF, il pezzo di filo e l'etichetta UHF.EN\_\allowbreak{}RX\_\allowbreak{}LNA, che non arrivavano a nessun pin (ERC label\_\allowbreak{}dangling) & Claude Code & \textbf{Brik81} & \textquotedblleft{}ok ho risolto facendo A\textquotedblright{} \\\midrule[0.15pt]
1.23 & proposta da Claude Code e approvata 27/\allowbreak{}10/\allowbreak{}2026\newline{}eseguita da Claude Code il 27/\allowbreak{}10 & S-11: D202 (protezione ESD di SCL/\allowbreak{}SDA dell'I2C1) passa da ESD2CAN24DBZRQ1, TVS per CAN a 24 V, a Nexperia PESD5V0L2BT,215 a 5 V. Stesso footprint SOT-23 e stessi pin & Claude Code & \textbf{Brik81} & \textquotedblleft{}se va bene per la can va bene anche per la i2c?\textquotedblright{} → \textquotedblleft{}ok, metti quello che consigli tu\textquotedblright{} \\\midrule[0.15pt]
1.24 & problemi R-01/\allowbreak{}R-02 dalla revisione indipendente del 27/\allowbreak{}10/\allowbreak{}2026\newline{}corretti da Brik81 in KiCad il 27/\allowbreak{}10 (salvataggio 14:14), verificati da Claude Code il 27/\allowbreak{}10 & R-01: il codice JLC di U301 passa da STM32L562RET6 (LQFP64 standard, pinout diverso dal simbolo) a STM32L562RET6P, coerente con il simbolo RETxP. R-02: il codice JLC di IC1/\allowbreak{}IC5 passa da TLV9001UIDBVR (pinout \textquotedblleft{}U\textquotedblright{} invertito) a TLV9051IDBVR & Brik81 (modifica) · Claude Code (verifica) & \textbf{Brik81} & \textquotedblleft{}R-01 solved; R-02 solved; (controlla e conferma nella rev pdf in cui indichiamo rev per rev ogni modifica fatta)\textquotedblright{} \\\midrule[0.15pt]
1.25 & problemi R-03/\allowbreak{}R-11 della revisione indipendente del 27/\allowbreak{}10/\allowbreak{}2026\newline{}modifica eseguita da Claude Code in \code{UHF.\allowbreak{}kicad\_\allowbreak{}sch} il 30/\allowbreak{}10 & L242 da 22 nH (LQW2BAN22NJ00L, rev 1.11) a 1 µH Taiyo Yuden NRS2012T1R0NGJV, footprint L\_\allowbreak{}Taiyo-Yuden\_\allowbreak{}NR-20xx; C260 da 33 pF a 27 pF GQM1875C2E270FB12D. Rete d'uscita riallineata alla EVB-02 Qorvo. Eccezione approvata alle regole \textquotedblleft{}induttori RF solo Murata LQW\textquotedblright{} e \textquotedblleft{}footprint esistenti\textquotedblright{}, solo per L242 & Claude Code & \textbf{Brik81} & \textquotedblleft{}si procedi con questi 3 punti che mi hai proposto\textquotedblright{} (30/\allowbreak{}10), dopo \textquotedblleft{}vorrei riamenre coerente con worvo\textquotedblright{} (= Qorvo) e \textquotedblleft{}non è un problema che non sia su jlc\textquotedblright{} \\\midrule[0.15pt]
1.26 & problemi R-05 e R-22 della revisione indipendente del 27/\allowbreak{}10/\allowbreak{}2026\newline{}R-05 corretto da Brik81 in KiCad (salvataggio 30/\allowbreak{}10 13:23), R-22 eseguito da Claude Code il 30/\allowbreak{}10 & R-05: i campi JLCPCB Parts di R206/\allowbreak{}R207, R202 e C35 ora corrispondono al valore. R-22: il campo Datasheet di U701 punta al datasheet Diodes 74AUP2G14, cioè al componente montato & Brik81 (R-05) · Claude Code (R-22 e verifica) & \textbf{Brik81} & \textquotedblleft{}allora R-05 dovrebbero essere risolte tutte, fai tu R-22\textquotedblright{} (30/\allowbreak{}10) \\\midrule[0.15pt]
1.27 & gruppo B della revisione indipendente del 27/\allowbreak{}10/\allowbreak{}2026 (violazioni delle regole di progetto)\newline{}eseguito da Claude Code il 30/\allowbreak{}10 & Solo le correzioni che non richiedono scelte: footprint SMD per TP701–TP703; C261, C702 e C704 a Murata GJM con lo stesso valore e footprint; numero del pin schermo di J15 allineato al footprint (S1→SH). Restano da decidere: R-10 per i 100 pF/\allowbreak{}56 pF e R-09 per connettore e linguette & Claude Code & \textbf{Brik81} & \textquotedblleft{}Partiamo dal gruppo B\textquotedblright{} (30/\allowbreak{}10), in risposta all'elenco con le correzioni proposte \\\midrule[0.15pt]
1.28 & gruppo B della revisione indipendente del 27/\allowbreak{}10/\allowbreak{}2026, punti rimasti\newline{}eseguito da Claude Code il 30/\allowbreak{}10 & R-10: C206, C207, C208, C701, C212, C213 (100 pF) e C255, C230 (56 pF) passano a Murata GQM1875 0603 con lo stesso valore. Eccezione approvata alla regola \textquotedblleft{}footprint esistenti\textquotedblright{}, perché in 0402 la serie GJM arriva solo a 47 pF e GQM1555/\allowbreak{}GQM2195 sono a stock 0 su JLC. R-09: J15 resta com'è (codice C393939, linguette del guscio THT), per decisione di Brik81 & Claude Code & \textbf{Brik81} & \textquotedblleft{}NO IL CONNETTORE TIPOC va bene cois R-09 lascia cosi come sta\textquotedblright{} e, sulla lista degli 8 condensatori, \textquotedblleft{}si procedi\textquotedblright{} (30/\allowbreak{}10) \\\midrule[0.15pt]
1.29 & problema R-04 della revisione indipendente del 27/\allowbreak{}10/\allowbreak{}2026\newline{}decisione di Brik81 del 02/\allowbreak{}11, registrata da Claude Code & Nessuna modifica allo schema. Il limite della potenza in uscita (≤ 31 dBm all'antenna) e la sequenza di trasmissione sono requisiti del firmware; misura obbligatoria sul prototipo; R242/\allowbreak{}R243 restano come regolazione della VAPC. Aggiunta la raccolta delle note per il firmware (§41) & Claude Code & \textbf{Brik81} & \textquotedblleft{}si se l'hw è corretto andiamo di A e lasciamo scritto al firmwarsita\textquotedblright{} (02/\allowbreak{}11) \\\midrule[0.15pt]
1.30 & problemi R-06 e R-24 della revisione indipendente del 27/\allowbreak{}10/\allowbreak{}2026\newline{}decisione di Brik81 del 02/\allowbreak{}11, registrata da Claude Code & Nessuna modifica. Requisito fissato da Brik81: il backup dell'RTC deve coprire solo pochi secondi senza alimentazione. Range −20…+60 °C del CPH3225A, assenza dal catalogo JLC e autonomia di circa 7 h accettati & Claude Code & \textbf{Brik81} & \textquotedblleft{}temperatura non ti preoccuapre, mancanza su jlcpcb non ti preoccupare, 7h non ti preoccupare mi serve solo come salvataggio di backup nel caso di alimentazione assente quidni un lavoro di qualche secondo a stare larghi\textquotedblright{} (02/\allowbreak{}11) \\\midrule[0.15pt]
1.31 & problema R-12 della revisione indipendente del 27/\allowbreak{}10/\allowbreak{}2026\newline{}decisione di Brik81 del 02/\allowbreak{}11, nota eseguita da Claude Code & J708 segue il pinout standard del dock. +12V, +3V3\_\allowbreak{}CONN (3V3 generato a bordo nel foglio POWER) e RS485.A1/\allowbreak{}B1 non sono usati su questa scheda: restano con le loro etichette e senza collegamenti. Aggiunta una nota di testo accanto a J708 nel foglio IF. Nessun no-connect sui pin: in KiCad un no-connect su un pin che ha anche un'etichetta è un conflitto & Claude Code & \textbf{Brik81} & \textquotedblleft{}R-12 si perche il connettore è standard, ma in quiesto caso noi abbiamo che sulla board i 12V non sono usati, i 3v3\_\allowbreak{}conn non sono usati perche venogono prodottti internametne nel shieet power infatti i 5V sono connessi, idem per rs485 per ora non usati ma no problem, se vuoi lasciarli con la label ma mettere il connettore NC a me sta bene\textquotedblright{} (02/\allowbreak{}11) \\\midrule[0.15pt]
1.32 & problema R-07 della revisione indipendente del 27/\allowbreak{}10/\allowbreak{}2026\newline{}decisione di Brik81 del 02/\allowbreak{}11, registrata da Claude Code & Nessuna modifica. Y201 resta ECS-TXO-20CSMV4-320-AY-TR (±0,5 ppm, −30…+85 °C). L'alternativa Epson TG2016SMN (−40…+85 °C) è stata valutata e non adottata & Claude Code & \textbf{Brik81} & \textquotedblleft{}A, lascia cosi perfavore\textquotedblright{} (02/\allowbreak{}11) \\\midrule[0.15pt]
1.33 & problema R-16 della revisione indipendente del 27/\allowbreak{}10/\allowbreak{}2026\newline{}ST collegato da Brik81 in KiCad (salvataggio 02/\allowbreak{}11 12:47), verificato da Claude Code & ST (pin 9) di IC4 TPS2121 collegato a GND, come indica TI quando lo stato non serve. OV2 resta a GND (protezione da sovratensione sulla USB disattivata): a terra si usa solo un alimentatore da banco di qualità. La priorità assoluta al 5 V del dock (IN1) è confermata e invariata & Brik81 (modifica) · Claude Code (verifica) & \textbf{Brik81} & \textquotedblleft{}ST connesso a gnd, verifica, OV2 dotn worry il nostro alimentatore è buono (500 euro da banco)\textquotedblright{} (02/\allowbreak{}11) \\\midrule[0.15pt]
1.34 & problema R-19 della revisione indipendente del 27/\allowbreak{}10/\allowbreak{}2026\newline{}R204/\allowbreak{}R205 aggiunte da Brik81 in KiCad (salvataggio 02/\allowbreak{}11 13:11), verificate da Claude Code & IC205 resta TPS3825-33. Per poterlo escludere in laboratorio senza fili volanti: R204 0R (montata) in serie tra IC205 RESET e RF\_\allowbreak{}SDN; R205 0R (DNP) tra RF\_\allowbreak{}MR\_\allowbreak{}N (PC4) e RF\_\allowbreak{}SDN. Aggiunte la misura M-03 e le note firmware F-05 (aggiornata) e F-10 & Brik81 (modifica) · Claude Code (verifica e log) & \textbf{Brik81} & \textquotedblleft{}ok io lascerei cosi e speriamo…\textquotedblright{}, poi \textquotedblleft{}potrei prevedere un solder bridge… lo disissaldo e metto 0 ohm che cosi è conensso diretto al MCU gpio cosi comando via firmware il powerup\textquotedblright{} e \textquotedblleft{}guarda R205 e R204 le ho aggiunte ora\textquotedblright{} (02/\allowbreak{}11) \\\midrule[0.15pt]
1.35 & problema R-13 della revisione indipendente del 27/\allowbreak{}10/\allowbreak{}2026\newline{}decisione di Brik81 del 02/\allowbreak{}11, registrata da Claude Code & Nessuna modifica e nessuna nota sullo schema. I pull-up R206/\allowbreak{}R207 (4k7) di I2C1 restano collegati al 3V3 tramite il solder jumper JP201, aperto di serie: è lo standard di progetto & Claude Code & \textbf{Brik81} & \textquotedblleft{}Lascia come è senza nota, è nostro standard fare cosi, next\textquotedblright{} (02/\allowbreak{}11) \\\midrule[0.15pt]
1.36 & problema R-14 della revisione indipendente del 27/\allowbreak{}10/\allowbreak{}2026\newline{}decisione di Brik81 del 02/\allowbreak{}11, registrata da Claude Code & Nessuna modifica. C39 (10 µF 0402 10 V X5R, C315248) resta sull'ingresso 5 V del dock, accettando la perdita di capacità per polarizzazione DC & Claude Code & \textbf{Brik81} & \textquotedblleft{}lascia com'è, next\textquotedblright{} (02/\allowbreak{}11) \\
\end{xltabular}}

{\small Copie dei file prima di ogni revisione: v1.0 in \code{scratchpad/\allowbreak{}backup/\allowbreak{}}, v1.1 in \code{scratchpad/\allowbreak{}backup\_\allowbreak{}v1/\allowbreak{}}, v1.2 in \code{scratchpad/\allowbreak{}backup\_\allowbreak{}v11/\allowbreak{}} (file come salvati da KiCad alle 16:37), v1.3 in \code{scratchpad/\allowbreak{}backup\_\allowbreak{}v12/\allowbreak{}} (salvataggio KiCad delle 17:11) (cartella temporanea della sessione Claude Code). Rev. 1.13: \code{UHF.\allowbreak{}kicad\_\allowbreak{}sch} e \code{SYMPHONY.\allowbreak{}kicad\_\allowbreak{}pro} prima della modifica salvati in \code{bak\_\allowbreak{}pre113/\allowbreak{}}. Rev. 1.12: \code{UHF.\allowbreak{}kicad\_\allowbreak{}sch} e \code{SYMPHONY.\allowbreak{}kicad\_\allowbreak{}pro} prima della modifica salvati in \code{bak\_\allowbreak{}pre112/\allowbreak{}}. Rev. 1.11: \code{UHF.\allowbreak{}kicad\_\allowbreak{}sch} prima della modifica salvato come \code{UHF.\allowbreak{}kicad\_\allowbreak{}sch.\allowbreak{}bak\_\allowbreak{}pre111}. Rev. 1.10: tutti i \code{*.\allowbreak{}kicad\_\allowbreak{}sch}, \code{SYMPHONY.\allowbreak{}kicad\_\allowbreak{}pro} e \code{UHF\_\allowbreak{}LIB/\allowbreak{}SYMBOL/\allowbreak{}*.\allowbreak{}kicad\_\allowbreak{}sym} prima della modifica salvati in \code{bak\_\allowbreak{}pre110/\allowbreak{}} nella cartella temporanea della sessione. Rev. 1.9: \code{UHF.\allowbreak{}kicad\_\allowbreak{}sch} prima della modifica salvato come \code{UHF.\allowbreak{}kicad\_\allowbreak{}sch.\allowbreak{}bak\_\allowbreak{}pre19} nella cartella temporanea della sessione. Rev. 1.8: \code{SYMPHONY.\allowbreak{}kicad\_\allowbreak{}sch} e \code{IF.\allowbreak{}kicad\_\allowbreak{}sch} prima della modifica salvati come \code{*.\allowbreak{}bak\_\allowbreak{}pre18} nella cartella temporanea della sessione. Rev. 1.7: \code{MCU.\allowbreak{}kicad\_\allowbreak{}sch} prima della modifica salvato come \code{MCU.\allowbreak{}kicad\_\allowbreak{}sch.\allowbreak{}bak\_\allowbreak{}pre17} nella cartella temporanea della sessione. Rev. 1.6: \code{MCU.\allowbreak{}kicad\_\allowbreak{}sch} e \code{SYMPHONY.\allowbreak{}kicad\_\allowbreak{}pcb} prima della modifica salvati come \code{*.\allowbreak{}bak\_\allowbreak{}pre16} nella cartella temporanea della sessione del 21/\allowbreak{}10. Rev. 1.5: \code{UHF.\allowbreak{}kicad\_\allowbreak{}sch} e \code{MCU.\allowbreak{}kicad\_\allowbreak{}sch} prima della modifica in \code{scratchpad/\allowbreak{}c2010/\allowbreak{}backup\_\allowbreak{}pre15/\allowbreak{}} (sessione del 20/\allowbreak{}10). Rev. 1.4: nessun file di progetto modificato; netlist ed ERC della verifica nella cartella temporanea della sessione del 20/\allowbreak{}10. Nella cartella del progetto c'è un solo PDF, sempre aggiornato all'ultima revisione: la storia delle modifiche sta nei registri qui sotto. Per uno storico permanente conviene fare un commit git per ogni revisione.}

\subsection[Registro modifiche]{Registro modifiche}

{\scriptsize\arrayrulecolor{rulec}\setlength{\tabcolsep}{2.5pt}\renewcommand{\arraystretch}{1.1}
\begin{xltabular}{\linewidth}{>{\hsize=0.4420\hsize\raggedright\arraybackslash}X>{\hsize=1.1366\hsize\raggedright\arraybackslash}X>{\hsize=1.2629\hsize\raggedright\arraybackslash}X>{\hsize=1.9575\hsize\raggedright\arraybackslash}X>{\hsize=1.3936\hsize\raggedright\arraybackslash}X>{\hsize=0.4181\hsize\raggedright\arraybackslash}X>{\hsize=0.6946\hsize\raggedright\arraybackslash}X>{\hsize=0.6946\hsize\raggedright\arraybackslash}X}
\toprule
\textbf{ID} & \textbf{Foglio} & \textbf{Ref} & \textbf{Prima} & \textbf{Dopo} & \textbf{Motivo (§)} & \textbf{Approvato da · quando} & \textbf{Eseguito · quando} \\
\midrule
\endhead
\midrule[0.15pt]\multicolumn{8}{r}{\footnotesize\itshape continua sulla pagina seguente\ldots}\\\bottomrule\endfoot
\bottomrule\endlastfoot
1.0-01 & UHF & Y201 /\allowbreak{} IC204 & TCXO su XOUT (pin 16), no-connect su XIN & TCXO su XIN (pin 17), no-connect su XOUT & §2 & Brik81 · 12:34 & Claude Code · 12:37 \\\midrule[0.15pt]
1.0-02 & UHF & C702 & 2p & 9.1p & §3 & Brik81 · 12:34 (condizione \textquotedblleft{}sicuro al 100\%\textquotedblright{}) & Claude Code · 12:42 \\\midrule[0.15pt]
1.0-03 & UHF & L702 & 3.2n & 33n & §3 & Brik81 · 12:34 & Claude Code · 12:42 \\\midrule[0.15pt]
1.0-04 & UHF & L701 → C704 & induttore 3.2n verso GND & condensatore C704 15p 0402 verso GND & §3 & Brik81 · 12:34 & Claude Code · 12:42 \\\midrule[0.15pt]
1.0-05 & UHF & C701 & 33p & 100p & §3 & Brik81 · 12:34 & Claude Code · 12:42 \\\midrule[0.15pt]
1.0-06 & UHF & L703 & 68n, DNP, scollegato & 220n montato, da 3V3 al pin TX (fili, giunzione, simbolo \#PWR0928) & §3 & Brik81 · 12:34 & Claude Code · 12:42 \\\midrule[0.15pt]
1.0-07 & UHF & R703 & 100k & 0 & §5 & Brik81 · 12:34 & Claude Code · 12:42 \\\midrule[0.15pt]
1.0-08 & UHF & R242 /\allowbreak{} R243 & 100 /\allowbreak{} 820 & 33 /\allowbreak{} 240 & §5 & Brik81 · 12:34 & Claude Code · 12:42 \\\midrule[0.15pt]
1.0-09 & UHF & C216 & 8.2 & 8.2p & §M & Brik81 · 12:34 & Claude Code · 12:42 \\\midrule[0.15pt]
1.0-10 & UHF & L243 & 2.4n & 2.7n & §M & Brik81 · 12:34 & Claude Code · 12:42 \\\midrule[0.15pt]
1.0-11 & UHF & C260 & 27p 0402 & 27p 0603 & §M & Brik81 · 12:34 & Claude Code · 12:42 \\\midrule[0.15pt]
1.0-12 & UHF & L242 & 1u 1206 & 0.9u 1008 & §M & Brik81 · 12:34 & Claude Code · 12:42 \\\midrule[0.15pt]
1.0-13 & UHF & C245, C256 & 3.3u 0402 & 10u 0805 & §M & Brik81 · 12:34 & Claude Code · 12:42 \\\midrule[0.15pt]
1.0-14 & UHF & L202 & 22n & 120n & §M & Brik81 · 12:34 & Claude Code · 12:42 \\\midrule[0.15pt]
1.0-15 & Libreria & Footprint IC2 & EP senza via & 9 via termiche (3×3) & §M & Brik81 · 12:34 & Claude Code · 12:44 \\\midrule[0.15pt]
1.0-16 & UHF & Note & – & 2 note di testo (match TX, layout uscita PA) & §M & Brik81 · 12:34 & Claude Code · 12:43–12:45 \\\midrule[0.15pt]
1.0-17 & UHF & Campo jlc & codici presenti & codici rimossi da C702, L702, C701, C704, L703, R703, L202 (valore cambiato) & §3, §M & Brik81 · 12:34 & Claude Code · 12:42 \\\midrule[0.15pt]
1.1-01 & POWER & R29 & 10k (jlc C71617) & 22k (jlc C25547) & §1.3 & Brik81 · 15:49 & Claude Code · 15:52 \\\midrule[0.15pt]
1.1-02 & POWER & R31 & 22k (jlc C25547) & 10k (jlc C71617) & §1.3 & Brik81 · 15:49 & Claude Code · 15:52 \\\midrule[0.15pt]
1.1-03 & UHF & Q701 & simbolo Transistor\_\allowbreak{}FET:2N7002E (N-MOS), descrizione e datasheet del 2N7002 & simbolo Transistor\_\allowbreak{}FET:Q\_\allowbreak{}PMOS\_\allowbreak{}GSD, descrizione \textquotedblleft{}P-Channel MOSFET −20 V −1.5 A\textquotedblright{} e datasheet onsemi FDN308P. Valore, footprint, jlc e collegamenti invariati & §6 & Brik81 · 15:49 & Claude Code · 15:52 \\\midrule[0.15pt]
1.1-04 & UHF & IC202 VDD (\#PWR0208) & 3V3 & 5V0 (C211 100n resta sul pin come disaccoppiamento) & §4 & Brik81 · 15:49 & Claude Code · 15:52 \\\midrule[0.15pt]
1.2-01 & UHF & R244 & 0 Ω tra APC2 e APC1 & rimosso, sostituito da un filo (APC1 e APC2 collegati direttamente) & §9 & Brik81 · 17:04 (scelta delegata a Claude Code) & Claude Code · 17:05 \\\midrule[0.15pt]
1.2-02 & UHF & C253 & 10n tra APC1 e GND (con simbolo GND \#PWR0911 e filo verso il nodo) & rimosso; resta C252 10n sul nodo APC unico & §9 & Brik81 · 17:04 (scelta delegata a Claude Code) & Claude Code · 17:05 \\\midrule[0.15pt]
1.3-01 & UHF & C201 & jlc: C1525 & jlc: C77014 = muRata GRM155R71E104KE14D & §10 & Brik81 · 17:11 & Claude Code · 17:13 \\\midrule[0.15pt]
1.3-02 & UHF & C203 & jlc: C1525 & jlc: C77014 = muRata GRM155R71E104KE14D & §10 & Brik81 · 17:11 & Claude Code · 17:13 \\\midrule[0.15pt]
1.3-03 & UHF & C211 & jlc: C1525 & jlc: C77014 = muRata GRM155R71E104KE14D & §10 & Brik81 · 17:11 & Claude Code · 17:13 \\\midrule[0.15pt]
1.3-04 & UHF & C202 & jlc: C29266 & jlc: C77009 = muRata GRM155R61E105KA12D & §10 & Brik81 · 17:11 & Claude Code · 17:13 \\\midrule[0.15pt]
1.3-05 & UHF & C215 & jlc: C76999 & jlc: C77009 = muRata GRM155R61E105KA12D & §10 & Brik81 · 17:11 & Claude Code · 17:13 \\\midrule[0.15pt]
1.3-06 & UHF & C204 & jlc: C21585 & jlc: C77013 = muRata GRM155R71E103KA01D & §10 & Brik81 · 17:11 & Claude Code · 17:13 \\\midrule[0.15pt]
1.3-07 & UHF & C206 & jlc: C1864265 & jlc: C21366878 = muRata GRM1555C2A101GA01D & §10 & Brik81 · 17:11 & Claude Code · 17:13 \\\midrule[0.15pt]
1.3-08 & UHF & C208 & jlc: C1864265 & jlc: C21366878 = muRata GRM1555C2A101GA01D & §10 & Brik81 · 17:11 & Claude Code · 17:13 \\\midrule[0.15pt]
1.3-09 & UHF & C207 & jlc: C1864265 & jlc: C77177 = muRata GRM1555C1H101JA01D & §10 & Brik81 · 17:11 & Claude Code · 17:13 \\\midrule[0.15pt]
1.3-10 & UHF & C230 & jlc: C437548 & jlc: C913752 = muRata GRM21A5C2E560JW01D & §10 & Brik81 · 17:11 & Claude Code · 17:13 \\\midrule[0.15pt]
1.3-11 & UHF & C240 & jlc: (vuoto) & jlc: C77013 = muRata GRM155R71E103KA01D & §10 & Brik81 · 17:11 & Claude Code · 17:13 \\\midrule[0.15pt]
1.3-12 & UHF & C246 & jlc: (vuoto) & jlc: C77013 = muRata GRM155R71E103KA01D & §10 & Brik81 · 17:11 & Claude Code · 17:13 \\\midrule[0.15pt]
1.3-13 & UHF & C252 & jlc: (vuoto) & jlc: C77013 = muRata GRM155R71E103KA01D & §10 & Brik81 · 17:11 & Claude Code · 17:13 \\\midrule[0.15pt]
1.3-14 & UHF & C241 & jlc: (vuoto) & jlc: C76947 = muRata GRM1555C1H102JA01D & §10 & Brik81 · 17:11 & Claude Code · 17:13 \\\midrule[0.15pt]
1.3-15 & UHF & C247 & jlc: (vuoto) & jlc: C76947 = muRata GRM1555C1H102JA01D & §10 & Brik81 · 17:11 & Claude Code · 17:13 \\\midrule[0.15pt]
1.3-16 & UHF & C257 & jlc: (vuoto) & jlc: C76947 = muRata GRM1555C1H102JA01D & §10 & Brik81 · 17:11 & Claude Code · 17:13 \\\midrule[0.15pt]
1.3-17 & UHF & C244 & jlc: (vuoto) & jlc: C76965 = muRata GRM1555C1H331JA01D & §10 & Brik81 · 17:11 & Claude Code · 17:13 \\\midrule[0.15pt]
1.3-18 & UHF & C248 & jlc: (vuoto) & jlc: C76965 = muRata GRM1555C1H331JA01D & §10 & Brik81 · 17:11 & Claude Code · 17:13 \\\midrule[0.15pt]
1.3-19 & UHF & C245 & jlc: (vuoto) & jlc: C84416 = muRata GRM21BR61E106KA73L & §10 & Brik81 · 17:11 & Claude Code · 17:13 \\\midrule[0.15pt]
1.3-20 & UHF & C256 & jlc: (vuoto) & jlc: C84416 = muRata GRM21BR61E106KA73L & §10 & Brik81 · 17:11 & Claude Code · 17:13 \\\midrule[0.15pt]
1.3-21 & UHF & C251 & jlc: (vuoto) & jlc: C77177 = muRata GRM1555C1H101JA01D & §10 & Brik81 · 17:11 & Claude Code · 17:13 \\\midrule[0.15pt]
1.3-22 & UHF & C258 & jlc: (vuoto) & jlc: C76970 = muRata GRM1555C1H470JA01D & §10 & Brik81 · 17:11 & Claude Code · 17:13 \\\midrule[0.15pt]
1.3-23 & UHF & C255 & jlc: (vuoto) & jlc: C393672 = muRata GCM1555C1H560JA16D & §10 & Brik81 · 17:11 & Claude Code · 17:13 \\\midrule[0.15pt]
1.3-24 & UHF & C260 & jlc: (vuoto) & jlc: C2176861 = muRata GCM1885C2A270FA16D & §10 & Brik81 · 17:11 & Claude Code · 17:13 \\\midrule[0.15pt]
1.3-25 & UHF & C261 & jlc: (vuoto) & jlc: C710855 = muRata GRM1555C2A220JA01D & §10 & Brik81 · 17:11 & Claude Code · 17:13 \\\midrule[0.15pt]
1.3-26 & UHF & C262 & jlc: (vuoto) & jlc: C161316 = muRata GJM1555C1H2R0BB01D & §10 & Brik81 · 17:11 & Claude Code · 17:13 \\\midrule[0.15pt]
1.3-27 & UHF & C701 & jlc: (vuoto) & jlc: C77177 = muRata GRM1555C1H101JA01D & §10 & Brik81 · 17:11 & Claude Code · 17:13 \\\midrule[0.15pt]
1.3-28 & UHF & C702 & jlc: (vuoto) & jlc: C415411 = muRata GRM1555C1H9R1DA01D & §10 & Brik81 · 17:11 & Claude Code · 17:13 \\\midrule[0.15pt]
1.3-29 & UHF & C704 & jlc: (vuoto) & jlc: C76950 = muRata GRM1555C1H150JA01D & §10 & Brik81 · 17:11 & Claude Code · 17:13 \\\midrule[0.15pt]
1.3-30 & UHF & L202 & jlc: (vuoto) & jlc: C113123 = muRata LQW15ANR12J00D & §10 & Brik81 · 17:11 & Claude Code · 17:13 \\\midrule[0.15pt]
1.3-31 & UHF & L240 & jlc: (vuoto) & jlc: C98068 = muRata LQW18AN12NG00D & §10 & Brik81 · 17:11 & Claude Code · 17:13 \\\midrule[0.15pt]
1.3-32 & UHF & L241 & jlc: (vuoto) & jlc: C86136 = muRata LQW18AN6N8C00D & §10 & Brik81 · 17:11 & Claude Code · 17:13 \\\midrule[0.15pt]
1.3-33 & UHF & L243 & jlc: (vuoto) & jlc: C412269 = muRata LQW15AN2N7G80D & §10 & Brik81 · 17:11 & Claude Code · 17:13 \\\midrule[0.15pt]
1.3-34 & UHF & L702 & jlc: (vuoto) & jlc: C1330869 = muRata LQW2BAN33NG00L & §10 & Brik81 · 17:11 & Claude Code · 17:13 \\\midrule[0.15pt]
1.3-35 & UHF & L703 & jlc: (vuoto) & jlc: C2046586 = muRata LQW18ANR22J80D & §10 & Brik81 · 17:11 & Claude Code · 17:13 \\\midrule[0.15pt]
1.3-36 & UHF & R240 & jlc: (vuoto) & jlc: C17168 = UNI-ROYAL 0402WGF0000TCE & §10 & Brik81 · 17:11 & Claude Code · 17:13 \\\midrule[0.15pt]
1.3-37 & UHF & R703 & jlc: (vuoto) & jlc: C17168 = UNI-ROYAL 0402WGF0000TCE & §10 & Brik81 · 17:11 & Claude Code · 17:13 \\\midrule[0.15pt]
1.3-38 & UHF & R241 & jlc: (vuoto) & jlc: C21189 = UNI-ROYAL 0603WAF0000T5E & §10 & Brik81 · 17:11 & Claude Code · 17:13 \\\midrule[0.15pt]
1.3-39 & UHF & R242 & jlc: (vuoto) & jlc: C25105 = UNI-ROYAL 0402WGF330JTCE & §10 & Brik81 · 17:11 & Claude Code · 17:13 \\\midrule[0.15pt]
1.3-40 & UHF & R243 & jlc: (vuoto) & jlc: C114755 = YAGEO RC0402FR-07240RL & §10 & Brik81 · 17:11 & Claude Code · 17:13 \\\midrule[0.15pt]
1.3-41 & UHF & R245 & jlc: (vuoto) & jlc: C25076 = UNI-ROYAL 0402WGF1000TCE & §10 & Brik81 · 17:11 & Claude Code · 17:13 \\\midrule[0.15pt]
1.3-42 & UHF & R247 & jlc: (vuoto) & jlc: C25076 = UNI-ROYAL 0402WGF1000TCE & §10 & Brik81 · 17:11 & Claude Code · 17:13 \\\midrule[0.15pt]
1.3-43 & UHF & R246 & jlc: (vuoto) & jlc: C27592 = UNI-ROYAL 0603WAF680JT5E & §10 & Brik81 · 17:11 & Claude Code · 17:13 \\\midrule[0.15pt]
1.3-44 & UHF & R248 & jlc: (vuoto) & jlc: C38941 = UNI-ROYAL 0402WGF1800TCE & §10 & Brik81 · 17:11 & Claude Code · 17:13 \\\midrule[0.15pt]
1.3-45 & UHF & IC2 & jlc: (vuoto) & jlc: C52059849 = Qorvo QPA9510TR7 & §10 & Brik81 · 17:11 & Claude Code · 17:13 \\\midrule[0.15pt]
1.3-46 & UHF & L243 & Footprint: Inductor\_\allowbreak{}SMD:L\_\allowbreak{}0603\_\allowbreak{}1608Metric\_\allowbreak{}Pad1.05x0.95mm\_\allowbreak{}HandSolder & Footprint: Inductor\_\allowbreak{}SMD:L\_\allowbreak{}0402\_\allowbreak{}1005Metric\_\allowbreak{}Pad0.77x0.64mm\_\allowbreak{}HandSolder & §10 & Brik81 · 17:11 & Claude Code · 17:13 \\\midrule[0.15pt]
1.3-47 & UHF & L703 & Footprint: Inductor\_\allowbreak{}SMD:L\_\allowbreak{}0805\_\allowbreak{}2012Metric & Footprint: Inductor\_\allowbreak{}SMD:L\_\allowbreak{}0603\_\allowbreak{}1608Metric\_\allowbreak{}Pad1.05x0.95mm\_\allowbreak{}HandSolder & §10 & Brik81 · 17:11 & Claude Code · 17:13 \\\midrule[0.15pt]
1.3-47 & UHF & L703 & Note: Lchoke - AN648 Tab.2 434MHz, alimenta il pin TX open-drain (es. Coilcraft 0805CS-221XJ, 400mA) & Note: Lchoke - AN648 Tab.2 434MHz, alimenta il pin TX open-drain: Murata LQW18ANR22J80D (280mA) & §10 & Brik81 · 17:11 & Claude Code · 17:13 \\\midrule[0.15pt]
\rowcolor{extrow}1.3+-01 & UHF & C226 & 100p 0402 tra 3V3 e GND (avanzo vecchia rete 5V\_\allowbreak{}RF) & rimosso & P5 & non registrato & KiCad, fuori sessione · 19/\allowbreak{}10 17:13–17:26 \\\midrule[0.15pt]
\rowcolor{extrow}1.3+-02 & UHF & C227 & 100p 0402 tra 3V3 e GND (avanzo vecchia rete 5V\_\allowbreak{}RF) & rimosso & P5 & non registrato & KiCad, fuori sessione · 19/\allowbreak{}10 17:13–17:26 \\\midrule[0.15pt]
\rowcolor{extrow}1.3+-03 & UHF & C228 & 1n 0402 tra 3V3 e GND (avanzo vecchia rete 5V\_\allowbreak{}RF) & rimosso & P5 & non registrato & KiCad, fuori sessione · 19/\allowbreak{}10 17:13–17:26 \\\midrule[0.15pt]
\rowcolor{extrow}1.3+-04 & IF & J3 & Molex 504050-0891, connettore di debug, DNP & rimosso & – & non registrato & KiCad, fuori sessione · 19/\allowbreak{}10 17:13–17:26 \\\midrule[0.15pt]
\rowcolor{extrow}1.3+-05 & IF & R10 & 22 Ω DNP, SWCLK → J3.4 & rimosso & – & non registrato & KiCad, fuori sessione · 19/\allowbreak{}10 17:13–17:26 \\\midrule[0.15pt]
\rowcolor{extrow}1.3+-06 & IF & R14 & 22 Ω DNP, SWDIO → J3.5 & rimosso & – & non registrato & KiCad, fuori sessione · 19/\allowbreak{}10 17:13–17:26 \\\midrule[0.15pt]
\rowcolor{extrow}1.3+-07 & IF & R15 & 22 Ω DNP, SWO → J3.6 & rimosso & – & non registrato & KiCad, fuori sessione · 19/\allowbreak{}10 17:13–17:26 \\\midrule[0.15pt]
\rowcolor{extrow}1.4+-01 & IF & IC501 & TCAN3414 a specchio & simbolo corretto, pin al posto giusto & P8 & Brik81 (autore) & Brik81 in KiCad · 20/\allowbreak{}10 09:31–14:38 \\\midrule[0.15pt]
\rowcolor{extrow}1.4+-02 & MCU & U301 pin 30, 62 & VDD12 su 3V3 & non collegati (NC) & P9 & Brik81 (autore) & Brik81 in KiCad · 20/\allowbreak{}10 09:31–14:38 \\\midrule[0.15pt]
\rowcolor{extrow}1.4+-03 & MCU & C1, C2 & 10 nF & 10 pF & P10 & Brik81 (autore) & Brik81 in KiCad · 20/\allowbreak{}10 09:31–14:38 \\\midrule[0.15pt]
\rowcolor{extrow}1.4+-04 & UHF & R202 /\allowbreak{} R203 & 33 k /\allowbreak{} 10 k & 0 Ω /\allowbreak{} valore \textquotedblleft{}DNP\textquotedblright{} & P11 & Brik81 (autore) & Brik81 in KiCad · 20/\allowbreak{}10 09:31–14:38 \\\midrule[0.15pt]
\rowcolor{extrow}1.4+-05 & UHF & U701 pin 3 & flottante & GND & P12 & Brik81 (autore) & Brik81 in KiCad · 20/\allowbreak{}10 09:31–14:38 \\\midrule[0.15pt]
\rowcolor{extrow}1.4+-06 & UHF & C207 /\allowbreak{} C208 (IC202) & PA su RF1, LNA su RF2 & LNA su RF1, PA su RF2 & sch\_\allowbreak{}check §6.6 & Brik81 (autore) & Brik81 in KiCad · 20/\allowbreak{}10 09:31–14:38 \\\midrule[0.15pt]
\rowcolor{extrow}1.4+-07 & IF & R206, R207 & 10 k & 4k7 & – & Brik81 (autore) & Brik81 in KiCad · 20/\allowbreak{}10 09:31–14:38 \\\midrule[0.15pt]
\rowcolor{extrow}1.4+-08 & SENSORS & R4–R7 & valore \textquotedblleft{}R\textquotedblright{} & 3.16 k /\allowbreak{} 105 k /\allowbreak{} 63.4 k /\allowbreak{} 19.6 k & P14 & Brik81 (autore) & Brik81 in KiCad · 20/\allowbreak{}10 09:31–14:38 \\\midrule[0.15pt]
\rowcolor{extrow}1.4+-09 & SENSORS & IC5, R8–R10, R12, R\_\allowbreak{}RTD1, C312, C313 & – & aggiunti (secondo canale) & P14 & Brik81 (autore) & Brik81 in KiCad · 20/\allowbreak{}10 09:31–14:38 \\\midrule[0.15pt]
1.5-01 & UHF & IC204 pin 1 (SDN) & a GND (\#PWR0204) & rete RF\_\allowbreak{}SDN → IC205 pin 3 (RESET) & S-01, §17 & Brik81 · 20/\allowbreak{}10 17:02 & Claude Code · 20/\allowbreak{}10 17:07–17:12 \\\midrule[0.15pt]
1.5-02 & UHF & IC205 (nuovo) & – & TPS3825-33DBVR, IC:SOT95P280X145-5N, jlc C473344. VDD 3V3, GND & S-01, §17 & Brik81 · 20/\allowbreak{}10 17:02 & Claude Code · 20/\allowbreak{}10 17:07–17:12 \\\midrule[0.15pt]
1.5-03 & UHF & C263 (nuovo) & – & 100n 0402 su VDD di IC205, jlc C77014 & S-01, §17 & Brik81 · 20/\allowbreak{}10 17:02 & Claude Code · 20/\allowbreak{}10 17:07–17:12 \\\midrule[0.15pt]
1.5-04 & UHF & R208 (nuovo) & – & 10k 0402 pull-up di \textasciitilde{}MR a 3V3, jlc C71617 & S-01, §17 & Brik81 · 20/\allowbreak{}10 17:02 & Claude Code · 20/\allowbreak{}10 17:07–17:12 \\\midrule[0.15pt]
1.5-05 & UHF & R209, D203 (nuovi) & – & \textasciitilde{}RESET → R209 1k5 0402 (jlc C25867) → D203 LED rosso 0402 NCD0402R1 (jlc C130719) → GND & S-01, §17 & Brik81 · 20/\allowbreak{}10 17:02 & Claude Code · 20/\allowbreak{}10 17:07–17:12 \\\midrule[0.15pt]
1.5-06 & UHF /\allowbreak{} MCU & U301 pin 24 (PC4) & non collegato & global label RF\_\allowbreak{}MR\_\allowbreak{}N → IC205 pin 4 (\textasciitilde{}MR) & S-01, §17 & Brik81 · 20/\allowbreak{}10 17:02 & Claude Code · 20/\allowbreak{}10 17:07–17:12 \\\midrule[0.15pt]
1.6-01 & MCU & BT401 & Battery\_\allowbreak{}Cell MS621FE-FL11E, footprint Seiko MS621F, jlc C9900147719, + su U404 pin 6 (VBAT) & rimossa & S-02, §18 & Brik81 · 21/\allowbreak{}10 17:44–17:55 & Claude Code · 21/\allowbreak{}10 17:55–17:58 \\\midrule[0.15pt]
1.6-02 & MCU & C415 (nuovo) & – & C\_\allowbreak{}Polarized 0.33F 5.5V, CDA CBL5R5334CF-ZJ, Capacitor\_\allowbreak{}THT:CP\_\allowbreak{}Radial\_\allowbreak{}D13.0mm\_\allowbreak{}P5.00mm, jlc C33990384. + su VBAT (U404 pin 6), − su GND & S-02, §18 & Brik81 · 21/\allowbreak{}10 17:44–17:55 & Claude Code · 21/\allowbreak{}10 17:55–17:58 \\\midrule[0.15pt]
1.6-03 & MCU & D401 (nuovo) & – & BAT54WS-7-F (Diodes Inc.), Diode\_\allowbreak{}SMD:D\_\allowbreak{}SOD-323, jlc C124205. Anodo su 3V3, catodo su R412 & S-02, §18 & Brik81 · 21/\allowbreak{}10 17:44–17:55 & Claude Code · 21/\allowbreak{}10 17:55–17:58 \\\midrule[0.15pt]
1.6-04 & MCU & R412 (nuovo) & – & 1k 0402, Resistor\_\allowbreak{}SMD:R\_\allowbreak{}0402\_\allowbreak{}1005Metric, jlc C11702 (Basic). Tra catodo di D401 e VBAT & S-02, §18 & Brik81 · 21/\allowbreak{}10 17:44–17:55 & Claude Code · 21/\allowbreak{}10 17:55–17:58 \\\midrule[0.15pt]
1.7-01 & MCU & C415 (1.6) & C\_\allowbreak{}Polarized 0.33F 5.5V CBL5R5334CF-ZJ, Capacitor\_\allowbreak{}THT:CP\_\allowbreak{}Radial\_\allowbreak{}D13.0mm\_\allowbreak{}P5.00mm, jlc C33990384 & rimosso da Brik81 (non SMD) & S-02, §19 & Brik81 · 23/\allowbreak{}10 (simbolo C5 incollato + richiesta) & Claude Code · 23/\allowbreak{}10 17:18 \\\midrule[0.15pt]
1.7-02 & MCU & C5 → C415 & cap:CPH3225A incollato da Brik81 con riferimento C5, pin 2 scollegato & rinominato C415; + su RTC\_\allowbreak{}VBAT, − su GND & S-02, §19 & Brik81 · 23/\allowbreak{}10 (simbolo C5 incollato + richiesta) & Claude Code · 23/\allowbreak{}10 17:18 \\\midrule[0.15pt]
1.7-03 & MCU & C416…C424 (nuovi) & – & 9 × cap:CPH3225A (Seiko CPH3225A 11 mF 3.3 V, footprint cap:CPH3225A, Mouser 628-CPH3225A), in parallelo a C415 & S-02, §19 & Brik81 · 23/\allowbreak{}10 (simbolo C5 incollato + richiesta) & Claude Code · 23/\allowbreak{}10 17:18 \\\midrule[0.15pt]
1.7-04 & MCU & Rete RTC\_\allowbreak{}VBAT & Net-(U404-VBAT): U404 pin 6 + C415(1.6) + R412 pin 1 & etichetta RTC\_\allowbreak{}VBAT su U404 pin 6, sulle due file di condensatori e su R412 pin 1. \#PWR0426 e \#PWR0427 (GND) sulle due file & S-02, §19 & Brik81 · 23/\allowbreak{}10 (simbolo C5 incollato + richiesta) & Claude Code · 23/\allowbreak{}10 17:18 \\\midrule[0.15pt]
1.7-05 & MCU & Nota di calcolo (testo) & – & Testo nel foglio MCU, inserito da Brik81 in KiCad (salvato 23/\allowbreak{}10 17:34): perché 10 × CPH3225A, Ctot = 110 mF (88 mF a −20 \%), V0 3.15 V (peggiore 3.0 V), Vmin 2.3 V, t = Ctot·(V0−Vmin)/\allowbreak{}I → 12.4 h tipico, 7.9 h a I max, 4.9 h caso peggiore, 3.9 h con 8 pezzi; carica τ = 110 s & S-02, §19 & Brik81 · 23/\allowbreak{}10 (eseguita da Brik81) & Brik81 · 23/\allowbreak{}10 17:34 \\\midrule[0.15pt]
1.8-01 & MCU & U301 PA0 /\allowbreak{} PA1 → PA9 /\allowbreak{} PA10 & UART4.TX su U301 pin 14 (PA0), UART4.RX su pin 15 (PA1), etichetta gerarchica UART4\{UART\} & UART1.TX su U301 pin 42 (PA9), UART1.RX su pin 43 (PA10), etichetta gerarchica UART1\{UART\} & §20 & Brik81 · 23/\allowbreak{}10 & Brik81 (in KiCad, prima del 23/\allowbreak{}10 17:1x) \\\midrule[0.15pt]
1.8-02 & Principale & Pin dei fogli MCU e IF & UART4\{UART\} (MCU) e UART4\{UART\} (IF) & UART1\{UART\} su entrambi & §20 & Brik81 · 23/\allowbreak{}10 & Claude Code · 23/\allowbreak{}10 \\\midrule[0.15pt]
1.8-03 & IF & Etichetta gerarchica + 4 etichette & UART4\{UART\}, 2 × UART4.TX, 2 × UART4.RX & UART1\{UART\}, 2 × UART1.TX, 2 × UART1.RX & §20 & Brik81 · 23/\allowbreak{}10 & Claude Code · 23/\allowbreak{}10 \\\midrule[0.15pt]
1.9-01 & UHF & C219 → CR2 & 7.5p, shunt da RXP a GND, jlc C3831627 & \textbf{2.2p} C0G 0805 (0805CG2R2C500NT, jlc C1801), in serie dall'uscita del filtro passa-basso (nodo C218/\allowbreak{}L204) a RXP. Footprint 0805 invariato & S-03, §21 & Brik81 · 23/\allowbreak{}10 & Claude Code · 23/\allowbreak{}10 \\\midrule[0.15pt]
1.9-02 & UHF & L205 → LR2 & 18n 0805, in serie verso il ramo RXP, jlc C2042168 & \textbf{56n} ±2\% avvolto in filo, Murata LQW18AN56NG00D, jlc C98079, footprint Inductor\_\allowbreak{}SMD:L\_\allowbreak{}0603\_\allowbreak{}1608Metric\_\allowbreak{}Pad1.05x0.95mm\_\allowbreak{}HandSolder; dall'uscita del filtro a RXN & S-03, §21 & Brik81 · 23/\allowbreak{}10 & Claude Code · 23/\allowbreak{}10 \\\midrule[0.15pt]
1.9-03 & UHF & C220 → CR1 & 7.5p, in serie verso il ramo RXN, jlc C3831627 & \textbf{4.7p} C0G 0805 (0805CG4R7C500NT, jlc C1820), da RXN a GND (nella posizione grafica dell'ex L206) & S-03, §21 & Brik81 · 23/\allowbreak{}10 & Claude Code · 23/\allowbreak{}10 \\\midrule[0.15pt]
1.9-04 & UHF & L206 → LR1 & 22n 0805, shunt da RXN a GND, jlc C2042193 & \textbf{62n} ±2\% avvolto in filo, Murata LQW18AN62NG00D, jlc C162649, footprint Inductor\_\allowbreak{}SMD:L\_\allowbreak{}0603\_\allowbreak{}1608Metric\_\allowbreak{}Pad1.05x0.95mm\_\allowbreak{}HandSolder; tra RXP e RXN & S-03, §21 & Brik81 · 23/\allowbreak{}10 & Claude Code · 23/\allowbreak{}10 \\\midrule[0.15pt]
1.9-05 & UHF & C222, C223 & 100p di blocco DC su RXP /\allowbreak{} RXN, jlc C1864265 & eliminati e sostituiti da filo: il blocco DC lo fa CR2 (C219) & S-03, §21 & Brik81 · 23/\allowbreak{}10 & Claude Code · 23/\allowbreak{}10 \\\midrule[0.15pt]
1.9-06 & UHF & GND \#PWR0225 e testo & GND di C219; testo \textquotedblleft{}latice balun\textquotedblright{} & GND eliminato; testo sostituito da \textquotedblleft{}RX match AN643 4 elementi (434 MHz)\textquotedblright{} & S-03, §21 & Brik81 · 23/\allowbreak{}10 & Claude Code · 23/\allowbreak{}10 \\\midrule[0.15pt]
1.9-07 & UHF & C219 (CR2) & 2.2p Fenghua 0805CG2R2C500NT, C0G generico, ±0.25 pF, jlc C1801, 0805 & 2.2p \textbf{Murata GQM1875C2E2R2BB12D} (alto Q per RF, ±0.1 pF), jlc \textbf{C2179335}, footprint Capacitor\_\allowbreak{}SMD:C\_\allowbreak{}0603\_\allowbreak{}1608Metric\_\allowbreak{}Pad1.08x0.95mm\_\allowbreak{}HandSolder. In 0805 tutti i GQM2195 da 2.2 pF hanno stock 0 & S-03, §21 & Brik81 · 24/\allowbreak{}10 & Claude Code · 24/\allowbreak{}10 \\\midrule[0.15pt]
1.9-08 & UHF & C220 (CR1) & 4.7p Fenghua 0805CG4R7C500NT, C0G generico, jlc C1820 & 4.7p \textbf{Murata GQM2195C2E4R7BB12D} (alto Q per RF, ±0.1 pF), jlc \textbf{C5646702}, 0805 invariato & S-03, §21 & Brik81 · 24/\allowbreak{}10 & Claude Code · 24/\allowbreak{}10 \\\midrule[0.15pt]
1.9-09 & UHF, MCU & C219, C220, L205, L206, D401, R412 & solo il campo \textquotedblleft{}jlc\textquotedblright{} & aggiunto il campo nascosto \textbf{\textquotedblleft{}JLCPCB PART\textquotedblright{}} con lo stesso codice di \textquotedblleft{}jlc\textquotedblright{} (C2179335, C5646702, C98079, C162649, C124205, C11702). Verificato nella BOM esportata con kicad-cli & S-03, §21 & Brik81 · 24/\allowbreak{}10 & Claude Code · 24/\allowbreak{}10 \\\midrule[0.15pt]
1.9-10 & UHF & IC203 TQP7M9106 e passivi & IC203, C224, C225, C229, C703, R204, R205, R701, L207, L208, L209, FB201: blocco tutto scollegato (32 pin\_\allowbreak{}not\_\allowbreak{}connected nell'ERC) & rimossi da Brik81 in KiCad (foglio salvato il 24/\allowbreak{}10 alle 15:53). Rilevato da Claude Code confrontando il foglio e confermato da Brik81: \textquotedblleft{}si l'ho tolto tanto era il blocco vecchio non usato ed era in DNP mi creava disordine\textquotedblright{}. Netlist dei componenti collegati invariata & S-03, §21 & Brik81 · 24/\allowbreak{}10 & Brik81 · 24/\allowbreak{}10 15:53 \\\midrule[0.15pt]
1.10-01 & tutti & Campo codice JLCPCB su 150 componenti & CONNECTORS 7, IF 18, MCU 32, POWER 22, UHF 86 campi, con nomi jlc /\allowbreak{} jlcpcb /\allowbreak{} JLCPCB Part /\allowbreak{} JLCPCB PART /\allowbreak{} JLCPCB & un solo campo \textbf{\textquotedblleft{}JLCPCB Parts\textquotedblright{}} per componente, valore invariato. Eliminati i doppioni con lo stesso valore: D401, R412 (MCU), IC205, C219, C220, L205, L206 (UHF). MCU 30 e UHF 81 campi, gli altri fogli invariati & §22 & Brik81 · 24/\allowbreak{}10 & Claude Code · 24/\allowbreak{}10 \\\midrule[0.15pt]
1.10-02 & librerie & UHF\_\allowbreak{}LIB/\allowbreak{}SYMBOL: IC, MASTER\_\allowbreak{}Dock\_\allowbreak{}SymLib, Power & JLCPCB PART, JLCPCB Part (×3), JLCPCB & \textquotedblleft{}JLCPCB Parts\textquotedblright{}, così i nuovi componenti piazzati da queste librerie nascono già col nome giusto & §22 & Brik81 · 24/\allowbreak{}10 & Claude Code · 24/\allowbreak{}10 \\\midrule[0.15pt]
1.10-03 & progetto & SYMPHONY.kicad\_\allowbreak{}pro, campi BOM & jlc, jlcpcb (mostrato e raggruppato), JLCPCB Part, JLCPCB & un solo campo \textquotedblleft{}JLCPCB Parts\textquotedblright{}, mostrato e raggruppato come il vecchio \textquotedblleft{}jlcpcb\textquotedblright{} & §22 & Brik81 · 24/\allowbreak{}10 & Claude Code · 24/\allowbreak{}10 \\\midrule[0.15pt]
1.10-04 & MCU, UHF & Formato file & 6 righe con a capo incompleto (solo CR), introdotte il 24/\allowbreak{}10 dallo script che aveva aggiunto \textquotedblleft{}JLCPCB PART\textquotedblright{} (riga 1.9-09) & corrette in CR+LF. Lo script è stato corretto & §22 & Brik81 · 24/\allowbreak{}10 & Claude Code · 24/\allowbreak{}10 \\\midrule[0.15pt]
1.11-01 & UHF & L242 & 0.9u, Inductor\_\allowbreak{}SMD:L\_\allowbreak{}1008\_\allowbreak{}2520Metric\_\allowbreak{}Pad1.43x2.20mm\_\allowbreak{}HandSolder, nessun codice JLCPCB & \textbf{22n} Murata \textbf{LQW2BAN22NJ00L} (RF avvolto in filo, nucleo ceramico, 1.9 A, 70 mΩ, ±5 \%), footprint Inductor\_\allowbreak{}SMD:L\_\allowbreak{}0805\_\allowbreak{}2012Metric\_\allowbreak{}Pad1.05x1.20mm\_\allowbreak{}HandSolder, JLCPCB Parts \textbf{C2042193} (stesso componente di L201/\allowbreak{}L203/\allowbreak{}L204) & §23 & Brik81 · 24/\allowbreak{}10 & Claude Code · 24/\allowbreak{}10 \\\midrule[0.15pt]
1.11-02 & UHF & C260 & 27p 0603, JLCPCB Parts C2176861 & \textbf{33p} Murata \textbf{GQM1875C2E330FB12D} (C0G alto Q, ±1 \%, 250 V), 0603 invariato, JLCPCB Parts \textbf{C2360428} & §23 & Brik81 · 24/\allowbreak{}10 & Claude Code · 24/\allowbreak{}10 \\\midrule[0.15pt]
1.12-01 & UHF & Reti RXP /\allowbreak{} RXN del Si4464 (IC204 pin 2 /\allowbreak{} 3) & nomi automatici Net-(IC204-RXP), Net-(IC204-RXN) & etichette \textbf{RF\_\allowbreak{}RX\_\allowbreak{}P} e \textbf{RF\_\allowbreak{}RX\_\allowbreak{}N} → reti /\allowbreak{}UHF/\allowbreak{}RF\_\allowbreak{}RX\_\allowbreak{}P, /\allowbreak{}UHF/\allowbreak{}RF\_\allowbreak{}RX\_\allowbreak{}N, riconosciute da KiCad come coppia differenziale. Collegamenti invariati & §24 & Brik81 · 24/\allowbreak{}10 & Claude Code · 24/\allowbreak{}10 \\\midrule[0.15pt]
1.12-02 & progetto & SYMPHONY.kicad\_\allowbreak{}pro, netclass & solo Default & 6 nuove netclass e 31 assegnazioni (tabella §24.2). Default invariata & §24 & Brik81 · 24/\allowbreak{}10 & Claude Code · 24/\allowbreak{}10 \\\midrule[0.15pt]
1.13-01 & UHF & 28 reti RF & nomi automatici Net-(…) & 28 etichette sui terminali dei componenti, nomi in §25.1. Collegamenti invariati & §25 & Brik81 · 24/\allowbreak{}10 & Claude Code · 24/\allowbreak{}10 \\\midrule[0.15pt]
1.13-02 & progetto & SYMPHONY.kicad\_\allowbreak{}pro, assegnazioni netclass & 28 assegnazioni sui nomi automatici & stesse assegnazioni sui nuovi nomi /\allowbreak{}UHF/\allowbreak{}… & §25 & Brik81 · 24/\allowbreak{}10 & Claude Code · 24/\allowbreak{}10 \\\midrule[0.15pt]
1.13-03 & UHF & Footprint di 27 componenti RF & …\_\allowbreak{}Pad…\_\allowbreak{}HandSolder & footprint standard della stessa misura (§25.2) & §25 & Brik81 · 24/\allowbreak{}10 & Claude Code · 24/\allowbreak{}10 \\\midrule[0.15pt]
1.15-01 & MCU & Y1, valore & \textquotedblleft{}Crystal\_\allowbreak{}GND24\textquotedblright{} (nome del simbolo) & \textbf{\textquotedblleft{}24MHz 10pF\textquotedblright{}}. Componente invariato: YXC X322524MMB4SI, 24 MHz, CL 10 pF, ±10 ppm, ESR 50 Ω, JLCPCB Parts C70583 & S-04, §27 & Brik81 · 24/\allowbreak{}10 & Claude Code · 24/\allowbreak{}10 \\\midrule[0.15pt]
1.15-02 & MCU & C1, C2, campo \textquotedblleft{}JLCPCB Parts\textquotedblright{} & \textquotedblleft{}C76946 \textquotedblright{} (con spazio finale) & \textquotedblleft{}C76946\textquotedblright{}. Nessun altro codice del progetto ha spazi & S-04, §27 & Brik81 · 24/\allowbreak{}10 & Claude Code · 24/\allowbreak{}10 \\\midrule[0.15pt]
1.16-01 & POWER & C39 (nuovo) & – & \textbf{10 µF 10 V X5R 0402}, Samsung CL05A106MP5NUNC, JLCPCB Parts \textbf{C315248}, tra rete 5V (IC4 pin 7 IN1) e GND, con \#PWR0934 & S-05, §28 & Brik81 · 24/\allowbreak{}10 & Claude Code · 24/\allowbreak{}10 \\\midrule[0.15pt]
1.16-02 & POWER & C40 (nuovo) & – & \textbf{100 nF 16 V X7R 0402}, Samsung CL05B104KO5NNNC, JLCPCB Parts \textbf{C1525} (già usato nel progetto), tra 5V e GND, con \#PWR0935 & S-05, §28 & Brik81 · 24/\allowbreak{}10 & Claude Code · 24/\allowbreak{}10 \\\midrule[0.15pt]
1.17-01 & SENSORS (CONNECTORS.kicad\_\allowbreak{}sch) & Etichette V\_\allowbreak{}TEMP\_\allowbreak{}PA, V\_\allowbreak{}TEMP\_\allowbreak{}BUCK & etichette locali & etichette gerarchiche di uscita, stessa posizione (uscita di IC1 e IC5 TLV9051) & S-06, §29 & Brik81 · 24/\allowbreak{}10 & Claude Code · 24/\allowbreak{}10 \\\midrule[0.15pt]
1.17-02 & MCU & U301 pin 26 PB1, pin 9 PC1 & non collegati & etichetta gerarchica V\_\allowbreak{}TEMP\_\allowbreak{}PA su PB1 (ADC1\_\allowbreak{}IN16), V\_\allowbreak{}TEMP\_\allowbreak{}BUCK su PC1 (ADC1\_\allowbreak{}IN2) & S-06, §29 & Brik81 · 24/\allowbreak{}10 & Claude Code · 24/\allowbreak{}10 \\\midrule[0.15pt]
1.17-03 & radice & Blocchi SENSORS e MCU & nessun pin per i sensori & pin di foglio V\_\allowbreak{}TEMP\_\allowbreak{}PA e V\_\allowbreak{}TEMP\_\allowbreak{}BUCK su entrambi i blocchi, collegati con fili che girano attorno al blocco IF & S-06, §29 & Brik81 · 24/\allowbreak{}10 & Claude Code · 24/\allowbreak{}10 \\\midrule[0.15pt]
1.18-01 & SENSORS & R4, R8 (3.16 k) & nessun codice & Yageo RT0402BRD073K16L, JLCPCB Parts \textbf{C852759} & S-06, §30 & Brik81 · 24/\allowbreak{}10 & Claude Code · 24/\allowbreak{}10 \\\midrule[0.15pt]
1.18-02 & SENSORS & R5, R12 (105 k) & nessun codice & Viking ARG02BTC1053, JLCPCB Parts \textbf{C2984407} (lo Yageo equivalente C852475 ha solo 5 pz) & S-06, §30 & Brik81 · 24/\allowbreak{}10 & Claude Code · 24/\allowbreak{}10 \\\midrule[0.15pt]
1.18-03 & SENSORS & R6, R9 (63.4 k) & nessun codice & Yageo RT0402BRD0763K4L, JLCPCB Parts \textbf{C852877} & S-06, §30 & Brik81 · 24/\allowbreak{}10 & Claude Code · 24/\allowbreak{}10 \\\midrule[0.15pt]
1.18-04 & SENSORS & R7, R10 (19.6 k) & nessun codice & Yageo RT0402BRD0719K6L, JLCPCB Parts \textbf{C852589} & S-06, §30 & Brik81 · 24/\allowbreak{}10 & Claude Code · 24/\allowbreak{}10 \\\midrule[0.15pt]
1.18-05 & SENSORS & R\_\allowbreak{}RTD1, R\_\allowbreak{}RTD2 (PT1000) & nessun codice; Description \textquotedblleft{}negative temperature coefficient\textquotedblright{} (sbagliata) & campi Manufacturer = YAGEO Nexensos, MPN = \textbf{32207638}; Description corretta; \textquotedblleft{}JLCPCB Parts\textquotedblright{} vuoto (acquisto separato, montaggio a mano) & S-06, §30 & Brik81 · 24/\allowbreak{}10 & Claude Code · 24/\allowbreak{}10 \\\midrule[0.15pt]
1.19-01 & UHF & R245 (100 Ω, attenuatore TX) & 0402, C25076 (62.5 mW: 49 mW dissipati, 78 \%) & \textbf{0603} Uni-Royal 0603WAF1000T5E 100 mW, JLCPCB Parts \textbf{C22775} (basic), footprint R\_\allowbreak{}0603\_\allowbreak{}1608Metric. Carico 49 \% & S-07, §31 & Brik81 · 24/\allowbreak{}10 & Claude Code · 24/\allowbreak{}10 \\\midrule[0.15pt]
1.19-02 & MCU & C307 (4.7 µF bulk VDD) & 0402, C23733: a 3.3 V resta circa il 30–40 \% della capacità & \textbf{0603} Samsung CL10A475KO8NNNC X5R \textbf{16 V}, JLCPCB Parts \textbf{C19666} (basic). A 3.3 V resta circa il 70–80 \% & S-14, §31 & Brik81 · 24/\allowbreak{}10 & Claude Code · 24/\allowbreak{}10 \\\midrule[0.15pt]
1.19-03 & UHF & L241 (6.8 nH, VCC2 del PA) & Murata LQW18AN6N8C00D 0603, 750 mA, C86136 & \textbf{0402} Murata \textbf{LQW15AN6N8G80D} (±2 \%, 1.45 A, 68 mΩ), JLCPCB Parts \textbf{C882467}. Stessa misura della scheda di valutazione Qorvo, e più corrente ammessa della 0603 precedente & S-14, §31 & Brik81 · 24/\allowbreak{}10 & Claude Code · 24/\allowbreak{}10 \\\midrule[0.15pt]
1.20-01 & MCU & U404 DS3231M pin 6 VBAT & scollegato: nelle netlist dalla rev. 1.8 alla 1.19 risultava \textquotedblleft{}unconnected-(U404-VBAT-Pad6)\textquotedblright{}, quindi i supercap C415…C424 di S-02 non arrivavano all'RTC & collegato alla rete RTC\_\allowbreak{}VBAT (con C415…C424 e R412). Rilevato da Claude Code nel controllo dei pin non collegati, corretto da Brik81 in KiCad & §32 & Brik81 · 24/\allowbreak{}10 17:50 & Brik81 · 24/\allowbreak{}10 17:50 \\\midrule[0.15pt]
1.20-02 & MCU & R410 (10 k, pull-up di 32KHZ) & montata & \textbf{DNP}, fuori BOM; footprint mantenuto sul PCB per poterla montare se servirà il 32 kHz & §32 & Brik81 · 24/\allowbreak{}10 & Claude Code · 24/\allowbreak{}10 \\\midrule[0.15pt]
1.20-03 & MCU & U301: PC3, PC6–PC15, PA0–PA4, PA15, PB0, PB2, PB4, PB5, PB8 (22 pin) & liberi senza simbolo, 22 avvisi ERC \textquotedblleft{}pin non collegato\textquotedblright{} & simbolo di non collegato (X) su ognuno. Restano usabili in futuro togliendo la X & §32 & Brik81 · 24/\allowbreak{}10 & Claude Code · 24/\allowbreak{}10 \\\midrule[0.15pt]
1.21-01 & libreria & UHF\_\allowbreak{}LIB/\allowbreak{}FOOTPRINT/\allowbreak{}passive\_\allowbreak{}v2.pretty/\allowbreak{}cap.pretty/\allowbreak{}CPH3225A.kicad\_\allowbreak{}mod & pad 1 a (−0.8, −0.3), pad 2 a (1.4, −0.3): origine spostata di 0.3 mm rispetto al centro; modello 3D con offset (0.3, 0.3, 0.59) & tutto il disegno traslato di (−0.3, +0.3) mm: pad a (∓1.1, 0), corpo, serigrafia, courtyard e testi centrati sull'origine. Offset del modello 3D portato a (0, 0, 0.59) per restare allineato. Percorso del file 3D (assoluto, cartella Download) invariato: lo sistema Brik81 & §33 & Brik81 · 24/\allowbreak{}10 & Claude Code · 24/\allowbreak{}10 \\\midrule[0.15pt]
1.21-02 & radice & Testo \textquotedblleft{}cd C:\textbackslash{}Users\textbackslash{}alex2\textbackslash{}…\textquotedblright{} incollato per sbaglio & presente & tolto & §33 & Brik81 · 24/\allowbreak{}10 & Brik81 · 24/\allowbreak{}10 17:50 \\\midrule[0.15pt]
1.21-03 & UHF & U701 74AUP2G14 & segnalato come possibile doppione & nessuna modifica: simbolo a più unità (due inverter + alimentazione), un solo componente in BOM & §33 & Brik81 · 24/\allowbreak{}10 & – \\\midrule[0.15pt]
1.22-01 & MCU & Bus UHF\{UHF\_\allowbreak{}SIG\_\allowbreak{}CTRL\}, uscita UHF.EN\_\allowbreak{}RX\_\allowbreak{}LNA (275.59, 30.48) & ingresso dal bus, filo ed etichetta senza nessun pin collegato nel foglio MCU & tolti. EN\_\allowbreak{}RX\_\allowbreak{}LNA resta generato in UHF da U701 (pin 6, inverso di EN\_\allowbreak{}TX\_\allowbreak{}ANT\_\allowbreak{}SW, comandato da PC2): la rete continua a collegare R202, TP703 e U701 pin 6 & §34 & Brik81 · 27/\allowbreak{}10 & Brik81 · 27/\allowbreak{}10 12:10 \\\midrule[0.15pt]
1.23-01 & IF & D202 (ESD su I2C1.SCL /\allowbreak{} I2C1.SDA) & TI ESD2CAN24DBZRQ1, JLCPCB Parts C5736151: lavoro 24 V, tensione lasciata passare circa 35–37 V, 3 pF & Nexperia \textbf{PESD5V0L2BT,215}, JLCPCB Parts \textbf{C12324}: 2 canali bidirezionali, lavoro 5 V, 75 pF. Footprint SOT-23-3 e pin invariati (1 SCL, 2 SDA, 3 GND). Aggiornati Datasheet e Description & S-11, §35 & Brik81 · 27/\allowbreak{}10 & Claude Code · 27/\allowbreak{}10 \\\midrule[0.15pt]
1.24-01 & MCU & U301 (campo JLCPCB Parts) & C1236672 = \textbf{STM32L562RET6} (LQFP64 senza SMPS, Fig.14 DS12736 p.86): pinout diverso dal simbolo dal pin 25 al 30 e dal 54 al 62 & C1236664 = \textbf{STM32L562RET6P} (LQFP64 con SMPS esterno, Fig.16 p.87), coerente con simbolo e netlist. Valore, footprint e collegamenti invariati. Stock LCSC al 27/\allowbreak{}10: 0 & R-01, §36 & Brik81 · 27/\allowbreak{}10 & Brik81 · 27/\allowbreak{}10 14:14 (verifica Claude Code) \\\midrule[0.15pt]
1.24-02 & SENSORS & IC1, IC5 (campo JLCPCB Parts) & C544301 = \textbf{TLV9001UIDBVR}: variante \textquotedblleft{}U\textquotedblright{} con pinout 1 IN+, 2 V−, 3 IN−, 4 OUT, 5 V+ (TLV9001 Fig.6-2 p.7) & C1849429 = \textbf{TLV9051IDBVR} (TI, SOT-23-5, stock 1879): pinout 1 OUT, 2 V−, 3 IN+, 4 IN−, 5 V+, uguale al simbolo (TLV9051 Tab.5-1 p.3) & R-02, §36 & Brik81 · 27/\allowbreak{}10 & Brik81 · 27/\allowbreak{}10 14:14 (verifica Claude Code) \\\midrule[0.15pt]
1.25-01 & UHF & L242 & 22n, L\_\allowbreak{}0805\_\allowbreak{}2012Metric, JLC C2042193 (Murata LQW2BAN22NJ00L, 1,9 A). Equivalente alla EVB solo a 435 MHz, insieme a C260 = 33 pF (rev 1.11, §23) & \textbf{1u}, \textbf{Inductor\_\allowbreak{}SMD:L\_\allowbreak{}Taiyo-Yuden\_\allowbreak{}NR-20xx}, Taiyo Yuden \textbf{NRS2012T1R0NGJV} (nuovo nome LCXND2020MKT1R0N0G, AEC-Q200): 1 µH ±30 \%, Isat 1,9 A, Itemp 1,7 A (ΔT 40 °C), DCR 84 mΩ max, 2,0×2,0×1,2 mm. Mouser 963-NRS2012T1R0NGJV; campo JLCPCB Parts vuoto (non a catalogo JLC). Collegamenti invariati (3V3 → RF\_\allowbreak{}PA\_\allowbreak{}OUT) & R-03, §37 & Brik81 · 30/\allowbreak{}10 & Claude Code · 30/\allowbreak{}10 \\\midrule[0.15pt]
1.25-02 & UHF & C260 & 33p, GQM1875C2E330FB12D (C2360428): valore di compensazione dei 22 nH & \textbf{27p}, Murata \textbf{GQM1875C2E270FB12D} (C2360423, 0603, C0G 1 \%, 250 V, stock LCSC 6690 al 30/\allowbreak{}10) = C10 della EVB-02. Footprint invariato & R-11, §37 & Brik81 · 30/\allowbreak{}10 & Claude Code · 30/\allowbreak{}10 \\\midrule[0.15pt]
1.26-01 & IF & R206, R207 (JLCPCB Parts) & C71617 = 10 kΩ, mentre il valore è 4k7 & C25900 = 4,7 kΩ 0402 (già usato su R30/\allowbreak{}R32) & R-05, §38 & Brik81 · 30/\allowbreak{}10 & Brik81 · 30/\allowbreak{}10 13:23 (verifica Claude Code) \\\midrule[0.15pt]
1.26-02 & UHF & R202 (JLCPCB Parts) & C25779 = 33 kΩ, mentre il valore è 0R (serie su SHUT\_\allowbreak{}DOWN dell'LNA) & C17168 = 0 Ω 0402 & R-05, §38 & Brik81 · 30/\allowbreak{}10 & Brik81 · 30/\allowbreak{}10 13:23 (verifica Claude Code) \\\midrule[0.15pt]
1.26-03 & POWER & C35 (JLCPCB Parts) & C1524 = 10 nF 50 V, mentre il valore è 100n & C1525 = 100 nF 16 V 0402 & R-05, §38 & Brik81 · 30/\allowbreak{}10 & Brik81 · 30/\allowbreak{}10 13:23 (verifica Claude Code) \\\midrule[0.15pt]
1.26-04 & UHF & U701 (Datasheet) e simbolo IC:74AUP2G14 in libreria & https:/\allowbreak{}/\allowbreak{}www.ti.com/\allowbreak{}lit/\allowbreak{}ds/\allowbreak{}symlink/\allowbreak{}sn74lvc2g14.pdf (SN74LVC2G14 TI) & https:/\allowbreak{}/\allowbreak{}www.diodes.com/\allowbreak{}assets/\allowbreak{}Datasheets/\allowbreak{}74AUP2G14.pdf. Il codice JLC C526337 è 74AUP2G14DW-7 Diodes SOT-363 (controllato su LCSC il 30/\allowbreak{}10). Valore, footprint e collegamenti invariati & R-22, §38 & Brik81 · 30/\allowbreak{}10 & Claude Code · 30/\allowbreak{}10 \\\midrule[0.15pt]
1.27-01 & UHF & TP701, TP702, TP703 & Connector\_\allowbreak{}Pin:Pin\_\allowbreak{}D0.7mm\_\allowbreak{}L6.5mm\_\allowbreak{}W1.8mm\_\allowbreak{}FlatFork (through-hole) & \textbf{TestPoint:TestPoint\_\allowbreak{}Pad\_\allowbreak{}D1.0mm} (pad SMD, come TP1). Reti invariate: RF\_\allowbreak{}XIN, EN\_\allowbreak{}TX\_\allowbreak{}ANT\_\allowbreak{}SW, EN\_\allowbreak{}RX\_\allowbreak{}LNA & R-08, §39 & Brik81 · 30/\allowbreak{}10 & Claude Code · 30/\allowbreak{}10 \\\midrule[0.15pt]
1.27-02 & UHF & C261 (22p 0402, adattamento PA) & C710855 GRM1555C2A220JA01D (GRM, ±5 \%) & C882567 \textbf{GJM1555C1H220FB01D} (C0G ±1 \%, 50 V, stock 47343) & R-10, §39 & Brik81 · 30/\allowbreak{}10 & Claude Code · 30/\allowbreak{}10 \\\midrule[0.15pt]
1.27-03 & UHF & C702 (9.1p 0402, C0 AN648) & C415411 GRM1555C1H9R1DA01D (±0,5 pF) & C497527 \textbf{GJM1555C1H9R1BB01D} (C0G ±0,1 pF, 50 V, stock 5366) & R-10, §39 & Brik81 · 30/\allowbreak{}10 & Claude Code · 30/\allowbreak{}10 \\\midrule[0.15pt]
1.27-04 & UHF & C704 (15p 0402, CM AN648) & C76950 GRM1555C1H150JA01D (±5 \%) & C441742 \textbf{GJM1555C1H150FB01D} (C0G ±1 \%, 50 V, stock 311272) & R-10, §39 & Brik81 · 30/\allowbreak{}10 & Claude Code · 30/\allowbreak{}10 \\\midrule[0.15pt]
1.27-05 & POWER & J15 USB-C, pin schermo & Il pin del simbolo si chiamava \textquotedblleft{}S1\textquotedblright{}, mentre nel footprint GCT\_\allowbreak{}USB4105 i 4 pad del guscio si chiamano \textquotedblleft{}SH\textquotedblright{}: sul PCB lo schermo non sarebbe stato collegato a GND & Pin rinominato in \textquotedblleft{}SH\textquotedblright{} (nella copia del simbolo dentro lo schema e nell'istanza), come nel simbolo standard KiCad 10 Connector:USB\_\allowbreak{}C\_\allowbreak{}Receptacle\_\allowbreak{}USB2.0\_\allowbreak{}16P. Rete: J15/\allowbreak{}SH su GND & R-09 (a), §39 & Brik81 · 30/\allowbreak{}10 & Claude Code · 30/\allowbreak{}10 \\\midrule[0.15pt]
1.28-01 & UHF & C206 (100 pF) & C21366878 GRM1555C2A101GA01D, 0402. RF\_\allowbreak{}SW\_\allowbreak{}COM → RF\_\allowbreak{}ANT, blocco DC verso l'antenna & C17308632 Murata \textbf{GQM1875C2E101FB12D (C0G ±1 \%, 250 V)}, \textbf{Capacitor\_\allowbreak{}SMD:C\_\allowbreak{}0603\_\allowbreak{}1608Metric}. Valore e collegamenti invariati & R-10, §40 & Brik81 · 30/\allowbreak{}10 & Claude Code · 30/\allowbreak{}10 \\\midrule[0.15pt]
1.28-02 & UHF & C207 (100 pF) & C77177 GRM1555C1H101JA01D, 0402. RF\_\allowbreak{}RX\_\allowbreak{}F1 → RF\_\allowbreak{}SW\_\allowbreak{}RX, blocco DC ramo RX & C17308632 Murata \textbf{GQM1875C2E101FB12D (C0G ±1 \%, 250 V)}, \textbf{Capacitor\_\allowbreak{}SMD:C\_\allowbreak{}0603\_\allowbreak{}1608Metric}. Valore e collegamenti invariati & R-10, §40 & Brik81 · 30/\allowbreak{}10 & Claude Code · 30/\allowbreak{}10 \\\midrule[0.15pt]
1.28-03 & UHF & C208 (100 pF) & C21366878 GRM1555C2A101GA01D, 0402. RF\_\allowbreak{}LPF\_\allowbreak{}OUT → RF\_\allowbreak{}SW\_\allowbreak{}TX, blocco DC ramo TX & C17308632 Murata \textbf{GQM1875C2E101FB12D (C0G ±1 \%, 250 V)}, \textbf{Capacitor\_\allowbreak{}SMD:C\_\allowbreak{}0603\_\allowbreak{}1608Metric}. Valore e collegamenti invariati & R-10, §40 & Brik81 · 30/\allowbreak{}10 & Claude Code · 30/\allowbreak{}10 \\\midrule[0.15pt]
1.28-04 & UHF & C701 (100 pF) & C77177 GRM1555C1H101JA01D, 0402. RF\_\allowbreak{}ATT\_\allowbreak{}IN → RF\_\allowbreak{}TX\_\allowbreak{}M2, blocco DC verso il pad & C17308632 Murata \textbf{GQM1875C2E101FB12D (C0G ±1 \%, 250 V)}, \textbf{Capacitor\_\allowbreak{}SMD:C\_\allowbreak{}0603\_\allowbreak{}1608Metric}. Valore e collegamenti invariati & R-10, §40 & Brik81 · 30/\allowbreak{}10 & Claude Code · 30/\allowbreak{}10 \\\midrule[0.15pt]
1.28-05 & UHF & C255 (56 pF) & C393672 GCM1555C1H560JA16D, 0402. RF\_\allowbreak{}ATT\_\allowbreak{}OUT → RF\_\allowbreak{}PA\_\allowbreak{}IN, C1 della EVB Qorvo & C237591 Murata \textbf{GQM1875C2E560JB12D (C0G ±5 \%, 250 V)}, \textbf{Capacitor\_\allowbreak{}SMD:C\_\allowbreak{}0603\_\allowbreak{}1608Metric}. Valore e collegamenti invariati & R-10, §40 & Brik81 · 30/\allowbreak{}10 & Claude Code · 30/\allowbreak{}10 \\\midrule[0.15pt]
1.28-06 & UHF & C212 (100 pF) & C1864265 GRM21A5C2E101FW01D, 0805. RF\_\allowbreak{}LNA\_\allowbreak{}IN → RF\_\allowbreak{}RX\_\allowbreak{}F2, blocco DC ingresso LNA & C17308632 Murata \textbf{GQM1875C2E101FB12D (C0G ±1 \%, 250 V)}, \textbf{Capacitor\_\allowbreak{}SMD:C\_\allowbreak{}0603\_\allowbreak{}1608Metric}. Valore e collegamenti invariati & R-10, §40 & Brik81 · 30/\allowbreak{}10 & Claude Code · 30/\allowbreak{}10 \\\midrule[0.15pt]
1.28-07 & UHF & C213 (100 pF) & C1864265 GRM21A5C2E101FW01D, 0805. RF\_\allowbreak{}RX\_\allowbreak{}F3 → RF\_\allowbreak{}LNA\_\allowbreak{}OUT, blocco DC uscita LNA & C17308632 Murata \textbf{GQM1875C2E101FB12D (C0G ±1 \%, 250 V)}, \textbf{Capacitor\_\allowbreak{}SMD:C\_\allowbreak{}0603\_\allowbreak{}1608Metric}. Valore e collegamenti invariati & R-10, §40 & Brik81 · 30/\allowbreak{}10 & Claude Code · 30/\allowbreak{}10 \\\midrule[0.15pt]
1.28-08 & UHF & C230 (56 pF) & C913752 GRM21A5C2E560JW01D, 0805. RF\_\allowbreak{}PA\_\allowbreak{}MATCH → RF\_\allowbreak{}LPF\_\allowbreak{}IN, uscita PA & C237591 Murata \textbf{GQM1875C2E560JB12D (C0G ±5 \%, 250 V)}, \textbf{Capacitor\_\allowbreak{}SMD:C\_\allowbreak{}0603\_\allowbreak{}1608Metric}. Valore e collegamenti invariati & R-10, §40 & Brik81 · 30/\allowbreak{}10 & Claude Code · 30/\allowbreak{}10 \\\midrule[0.15pt]
1.28-09 & POWER & J15 USB-C & C393939 (SHOU HAN TYPE-C16PIN) con footprint GCT\_\allowbreak{}USB4105; 4 linguette del guscio THT & Nessuna modifica, per decisione di Brik81. Il pin schermo era già stato corretto in SH nella rev 1.27 & R-09, §40 & Brik81 · 30/\allowbreak{}10 & — \\\midrule[0.15pt]
1.29-01 & UHF & Catena TX (IC204 → pad → IC2 → FL201 → IC202) & Con il Si4464 a +20 dBm il PA è in saturazione profonda (≈ +9 dBm in ingresso, ne bastano ≈ −9). Al commutatore arrivano ≈ 32,5 dBm, contro un massimo assoluto dell'HMC349 di 33,5 dBm (30 dBm in hot switching); l'LFCG-490+ a 85 °C regge ≈ 1,9 W & Nessuna modifica all'hardware. Requisiti firmware F-01 e F-02 (§41) e misura M-01 sul prototipo. Se Pout supera i 31 dBm anche con il limite firmware, si ritoccano R242/\allowbreak{}R243 & R-04, §41 & Brik81 · 02/\allowbreak{}11 & — \\\midrule[0.15pt]
1.30-01 & MCU & C415–C424 (10× CPH3225A), D401, R412, U404 & R-06: CPH3225A specificato −20…+60 °C, non a catalogo JLC. R-24: autonomia ≈ 7 h (110 mF da 3,0 a 2,3 V con ≈ 3 µA) & Nessuna modifica. Il requisito è di pochi secondi, quindi il margine è di molti ordini di grandezza (ne basterebbero già alcune decine di µF) & R-06, R-24, §42 & Brik81 · 02/\allowbreak{}11 & — \\\midrule[0.15pt]
1.31-01 & IF & J708 (FSI-110-03-G-D-AB), pin 2/\allowbreak{}4 +3V3\_\allowbreak{}CONN, 6/\allowbreak{}8 +12V, 16 RS485.B1, 18 RS485.A1 & Reti con un solo nodo, senza spiegazione sullo schema (la nota di Brik81 spiegava solo +3V3\_\allowbreak{}CONN) & Nessun cambio elettrico. Nuova nota di testo a (185,42; 116,84) mm: \textquotedblleft{}+12V, RS485.A1/\allowbreak{}B1: non usati su questa scheda\textquotedblright{}. Le etichette restano per documentare il pinout standard & R-12, §43 & Brik81 · 02/\allowbreak{}11 & Claude Code · 02/\allowbreak{}11 \\\midrule[0.15pt]
1.31-02 & IF & J708 pin 13 GPIOSX\_\allowbreak{}2, 14 GPIOSX\_\allowbreak{}0, 17 GPIO\_\allowbreak{}PPS & Etichette senza collegamento all'MCU & Nessuna modifica: restano solo per i nomi standard del connettore. R-12 chiuso & R-12, §43 & Brik81 · 02/\allowbreak{}11 & — \\\midrule[0.15pt]
1.32-01 & UHF & Y201 TCXO 32 MHz & ECS-TXO-20CSMV4-320-AY-TR (C19858587): A = ±0,5 ppm, Y = −30…+85 °C (datasheet ECS, guida alla numerazione, pag. 1) & Nessuna modifica, per decisione di Brik81 & R-07, §44 & Brik81 · 02/\allowbreak{}11 & — \\\midrule[0.15pt]
1.33-01 & POWER & IC4 TPS2121, pin 9 ST & Non collegato (rete unconnected-(IC4-ST-Pad9)) & Collegato alla linea GND con un filo e una giunzione, verso il pin GND 12 di IC4. Datasheet TI SLVSEA3F, Pin Functions: \textquotedblleft{}Status output… Connect to GND if not required\textquotedblright{} & R-16, §45 & Brik81 · 02/\allowbreak{}11 & Brik81 · 02/\allowbreak{}11 12:47 (verifica Claude Code) \\\midrule[0.15pt]
1.33-02 & POWER & IC4 TPS2121, pin 4 OV2 & A GND: supervisore di sovratensione su IN2 (VBUS USB) disattivato & Nessuna modifica, per decisione di Brik81: a terra la USB viene alimentata solo da un alimentatore da banco affidabile. Sul dock (IN1) la protezione resta attiva a 6,02 V (R29/\allowbreak{}R30) & R-16, §45 & Brik81 · 02/\allowbreak{}11 & — \\\midrule[0.15pt]
1.34-01 & UHF & R204 (nuova) & IC205 pin 3 (RESET, push-pull) collegato direttamente a RF\_\allowbreak{}SDN & R204 0R 0402 (C17168), montata, in serie: IC205 pin 3 → nuova rete Net-(IC205-RESET) → R204 → /\allowbreak{}UHF/\allowbreak{}RF\_\allowbreak{}SDN (IC204 pin 1) & R-19, §46 & Brik81 · 02/\allowbreak{}11 & Brik81 · 02/\allowbreak{}11 13:11 (verifica Claude Code) \\\midrule[0.15pt]
1.34-02 & UHF & R205 (nuova, DNP) & — & R205 0402 (C17168), DNP, tra RF\_\allowbreak{}MR\_\allowbreak{}N (PC4 /\allowbreak{} U301 pin 24, pull-up R208, IC205 \textasciitilde{}MR) e /\allowbreak{}UHF/\allowbreak{}RF\_\allowbreak{}SDN. Il campo Value oggi vale \textquotedblleft{}DNP\textquotedblright{}; consigliato \textquotedblleft{}0R\textquotedblright{} con il solo flag DNP & R-19, §46 & Brik81 · 02/\allowbreak{}11 & Brik81 · 02/\allowbreak{}11 13:11 (verifica Claude Code) \\\midrule[0.15pt]
1.35-01 & IF & R206, R207 (4k7), JP201 (SolderJumper\_\allowbreak{}2\_\allowbreak{}Open) & Con JP201 aperto il bus I2C1 (PB6/\allowbreak{}PB7, J203, J708 pin 9/\allowbreak{}11) non ha pull-up a bordo & Nessuna modifica. Standard di progetto: JP201 si chiude solo sulla scheda che deve fornire i pull-up al bus & R-13, §47 & Brik81 · 02/\allowbreak{}11 & — \\\midrule[0.15pt]
1.36-01 & POWER & C39 (10 µF, rete 5V, IN1 di IC4) & Samsung CL05A106MP5NUNC 0402 10 V X5R (C315248): con 5 V applicati la capacità effettiva è molto sotto il nominale (curva DC-bias non verificata) & Nessuna modifica, per decisione di Brik81. Alternativa scartata: stesso componente di C38 (C15850, 0805 25 V) & R-14, §48 & Brik81 · 02/\allowbreak{}11 & — \\
\end{xltabular}}

{\small Le righe 1.3+ non sono state eseguite da Claude Code: le ho ricostruite dal confronto delle netlist. Brik81 è pregato di confermare chi le ha fatte e quando.}

\subsection[Registro decisioni (nessuna modifica ai file)]{Registro decisioni (nessuna modifica ai file)}

{\footnotesize\arrayrulecolor{rulec}\setlength{\tabcolsep}{4pt}\renewcommand{\arraystretch}{1.1}
\begin{xltabular}{\linewidth}{>{\hsize=0.1820\hsize\raggedright\arraybackslash}X>{\hsize=0.6721\hsize\raggedright\arraybackslash}X>{\hsize=1.7097\hsize\raggedright\arraybackslash}X>{\hsize=0.8336\hsize\raggedright\arraybackslash}X>{\hsize=1.6027\hsize\raggedright\arraybackslash}X}
\toprule
\textbf{ID} & \textbf{Data e ora} & \textbf{Decisione} & \textbf{Deciso da} & \textbf{Esito della verifica} \\
\midrule
\endhead
\midrule[0.15pt]\multicolumn{5}{r}{\footnotesize\itshape continua sulla pagina seguente\ldots}\\\bottomrule\endfoot
\bottomrule\endlastfoot
D1 & 19/\allowbreak{}10/\allowbreak{}2026 16:37 & Power mux TPS2121: se dock 5 V (IN1) e Type-C (IN2) sono collegate insieme, \textbf{deve vincere la dock} (obbligatorio). L'abilitazione del 3V3 da EN\_\allowbreak{}DOCK resta com'è & Brik81, dopo confronto con il reparto (16:37 e 16:45) & Già soddisfatto dalla rev. 1.1: PR1 = 1.60 V > 1.06 V con la dock a 5 V; la dock perde la priorità sotto 3.32 V ed è esclusa sopra 6.02 V. OV2 e CP2 a GND = funzioni disattivate (datasheet TPS2121 §9.5, Tabella 9-3, §10.2.4.1) \\
\end{xltabular}}

\subsection[Proposte in attesa di approvazione]{Proposte in attesa di approvazione}

{\footnotesize\arrayrulecolor{rulec}\setlength{\tabcolsep}{4pt}\renewcommand{\arraystretch}{1.1}
\begin{xltabular}{\linewidth}{>{\hsize=0.2044\hsize\raggedright\arraybackslash}X>{\hsize=2.4312\hsize\raggedright\arraybackslash}X>{\hsize=0.5276\hsize\raggedright\arraybackslash}X>{\hsize=0.8368\hsize\raggedright\arraybackslash}X}
\toprule
\textbf{\#} & \textbf{Proposta} & \textbf{§} & \textbf{Stato} \\
\midrule
\endhead
\midrule[0.15pt]\multicolumn{4}{r}{\footnotesize\itshape continua sulla pagina seguente\ldots}\\\bottomrule\endfoot
\bottomrule\endlastfoot
P1 & SDN del Si4464 & – & \tg{tagdone}{chiusa in rev. 1.5 (supervisore, §17)} \\\midrule[0.15pt]
P2 & L242: scelta finale (Isat 1.4 A attuale, oppure 2 A con footprint dedicato). Nessuna alternativa Murata adatta su JLCPCB & §M, §10 & decisione di Brik81 \\\midrule[0.15pt]
P3 & Semplificazione passivi di IC2 & §9 & \textbf{eseguita in rev. 1.2} (R244 e C253 rimossi; R240/\allowbreak{}R241 tenuti per il prototipo) \\\midrule[0.15pt]
P4 & Correzione codici JLCPCB e assegnazione codici Murata & §10 & \textbf{eseguita in rev. 1.3}, tranne L242 (P2) e J201 (P7) \\\midrule[0.15pt]
P7 & J201: il footprint è il Molex 73415-0961 (edge-mount per PCB da 0.8 mm), il codice C2999837 è un DreamLNK DL-MMCX-KEF-180. Scegliere il connettore e allineare codice e footprint. Il Molex su JLCPCB è C588476 (stock 12, catalogato come \textquotedblleft{}pre-order\textquotedblright{}) & §10.2 & decisione di Brik81 \\\midrule[0.15pt]
P5 & C226, C227, C228: avanzi del disaccoppiamento della vecchia rete 5V\_\allowbreak{}RF (IC601 + vecchio PA IC203). Proposta: riusare C227 100p vicino ai pin VDD di IC204 e aggiungere un 1 µF (come Si4464 datasheet Fig. 1: 100p + 100n + 1µ); togliere C226 e C228 & – & \textbf{chiusa (1.3+)}: C226, C227 e C228 rimossi in KiCad. Il Si4464 ha comunque 100p (C214) + 100n (C201, C203) + 1µ (C202, C215), come in Fig. 1 (§13.5) \\\midrule[0.15pt]
P6 & R29 da 22k 5 \% a 22k 1 \%: la soglia OV peggiore sale da 5.46 a 5.64 V (solo codice JLC) & §1.3 & da approvare \\\midrule[0.15pt]
P8 & \tg{tagcrit}{critica} IF – IC501 TCAN3414 collegato a specchio: specchiare il simbolo in orizzontale (i fili esistenti cadono al posto giusto) e collegare SHDN (pin 5) a GND o a un GPIO & §13.3 & \tg{tagdone}{chiusa (1.4+-01)} \\\midrule[0.15pt]
P9 & \tg{tagcrit}{critica} MCU – pin 30 e 62 (VDD12) di U301 scollegati dal 3V3, con marker no-connect. Il componente montato deve essere la versione RET6P & §13.4 & \tg{tagdone}{chiusa (1.4+-02)} \\\midrule[0.15pt]
P10 & \tg{tagcrit}{critica} MCU – C1 e C2 da 10 nF a \textbf{10 pF} C0G; indicare frequenza e C\textsubscript{L} di Y1 & §13.4 & \tg{tagdone}{chiusa (1.4+-03)} \\\midrule[0.15pt]
P11 & \tg{tagwarn}{media} UHF – shutdown dell'LNA: R202 = 0 Ω (o 1 kΩ) e R203 DNP, così il pin arriva a circa 3.3 V e l'LNA si spegne in TX & §13.5 & \tg{tagdone}{chiusa (1.4+-04)} \\\midrule[0.15pt]
P12 & \tg{tagnote}{minore} UHF – U701 pin 3 (ingresso 2A, inutilizzato) a GND & §13.5 & \tg{tagdone}{chiusa (1.4+-05)} \\\midrule[0.15pt]
P13 & \tg{tagnote}{minore} POWER – 10 µF + 100 nF su IN1 di IC4 (5V dalla dock), vicino al pin 7 & §13.2 & da approvare (rev. 1.4) \\\midrule[0.15pt]
P14 & \tg{tagcrit}{da completare} SENSORS – valori di R4–R7 dalla simulazione (3.16k /\allowbreak{} 105k /\allowbreak{} 63.4k /\allowbreak{} 19.6k, 0.1 \%) e V\_\allowbreak{}TEMP\_\allowbreak{}PA su un pin ADC (PA2, PA3, PA4, PC1, PB0 o PB1); secondo sensore & §13.6 & \tg{tagdone}{chiusa (1.4+-08/\allowbreak{}09); collegamento ADC da verificare} \\\midrule[0.15pt]
P15 & \tg{tagnote}{decisione} IF – segnali della dock GPIOSX\_\allowbreak{}0, GPIOSX\_\allowbreak{}2, GPIO\_\allowbreak{}PPS, RS485.A1/\allowbreak{}B1 non collegati (RS485 senza transceiver); pull-up I2C1 attivi solo chiudendo JP201: chiudere di default o garantire i pull-up sulla dock & §13.3 & da approvare (rev. 1.4) \\\midrule[0.15pt]
P16 & \tg{tagnote}{pulizia} Etichette residue (UHF.EN\_\allowbreak{}RX\_\allowbreak{}LNA sulla radice, 5v\_\allowbreak{}TX\_\allowbreak{}EN su PC3, CAN1\{CAN\}); DNP avanzati in UHF (IC203 \& co., 47 avvisi ERC); riferimento R\_\allowbreak{}RTD; simbolo TPSM863252; LED D9 con 10 kΩ (debole) & §13.1–§13.6 & da approvare (rev. 1.4) \\\midrule[0.15pt]
P17 & \tg{tagnote}{minore} IF – D202 (ESD2CAN24, TVS per CAN da 24 V) sull'I2C: meglio un TVS da 3.3–5 V di lavoro & §13.3 & da approvare (rev. 1.4) \\
\end{xltabular}}
